\documentclass[11pt]{article}
\usepackage[T1]{fontenc}
\usepackage{lmodern}
\usepackage[margin=1in]{geometry}
\usepackage{amsmath,amssymb,amsthm,mathtools,microtype,booktabs}
\usepackage{xcolor}
\usepackage[colorlinks=true,linkcolor=blue!50!black,urlcolor=blue!50!black,citecolor=blue!50!black]{hyperref}
\hypersetup{pdftitle={Effective computation and boundary robustness of the random friend-tree exponent},pdfauthor={Christopher D. Long}}
\setlength{\emergencystretch}{2em}
\allowdisplaybreaks
\newtheorem{theorem}{Theorem}[section]
\newtheorem{proposition}[theorem]{Proposition}
\newtheorem{lemma}[theorem]{Lemma}
\newtheorem{corollary}[theorem]{Corollary}
\theoremstyle{definition}
\newtheorem{definition}[theorem]{Definition}
\theoremstyle{remark}
\newtheorem{remark}[theorem]{Remark}
\title{Effective computation and boundary robustness\\
of the random friend-tree exponent}
\author{Christopher D. Long}
\date{Research draft --- November 23, 2026}
\begin{document}
\maketitle
\begin{abstract}
The random friend tree selects a uniformly random existing vertex and
attaches a new leaf to a uniformly random neighbour of that vertex.
Using the established existence of its mean growth exponent $\mu$, we
derive an exact characteristic equation and prove effective two-sided
computability. Integrating all leaf-clock events gives an explicit
beta kernel. The discounted genealogical coefficients of an inert
newborn component have a genuine generation limit $r_\infty(s)$,
which is continuous, strictly decreasing, and satisfies
$r_\infty(\mu)=1$. A uniform estimate for the resolvent effect of one
boundary leaf identifies the exponent of the supremum over all newborn
environments with $\mu$. Finite-profile time-block calculations then
provide complete upper certificates, complementing the exact positive
lower certificates and giving a terminating approximation algorithm
at every rational tolerance. The same resolvent estimate proves that
every bounded predictable root-leaf input has mean exponent $\mu$,
uniformly over the input policy. Finally, a rational correction of
the existing motif potential improves the certified upper bound to
$0.5664485422$. The analysis does not classify the arithmetic nature
of the original exponent or assert an efficient high-precision solver.
\end{abstract}
\clearpage
\setcounter{tocdepth}{1}
\tableofcontents
\clearpage

\section{Introduction and main results}
\label{intro:section}

The random friend tree grows by selecting a uniformly random existing
vertex, selecting one of its neighbours uniformly, and attaching a new
leaf to that neighbour. Krapivsky and Redner introduced the model in
their study of emergent network modularity~\cite{KR}. Its number $C_n$
of non-leaf vertices grows on a nontrivial power scale. We study the
exponent $\mu$ characterized by
\[
 \mathbb E C_n=n^{\mu+o(1)}.
\]
The existence of this exponent and its identification with the
almost-sure logarithmic exponent are inputs from Uzay's
computer-assisted work~\cite{Uzay}. That work gives
$0.5650\le\mu\le0.5665$, excluding $17/30$.
The broader probabilistic structure of random friend trees is developed
in~\cite{AB}.

Our main questions concern exact characterization and effective
approximation. An infinite operator equation can identify a real number
without supplying a terminating algorithm for its approximation.
Likewise, numerical evidence of agreement between two initial
environments does not imply that their mean exponents agree uniformly
over all environments. We address both issues by controlling the
effect of a single boundary leaf on an arbitrary old component.

We embed the tree process in continuous time by assigning a rate-one
source clock to each real vertex. A \emph{core} is a vertex of full
degree at least two. In a retained component, the root has one inert
parent stub: it contributes one to the root's degree, supplies no
incoming leaves, and produces a null transition when selected. The
\emph{inert newborn} consists of a root core, one real leaf, and this
stub. Write $F_\infty(t)$ for its expected number of cores.

For a newborn core in a finite environment $x$, let $F_x(t)$ be the
expected number of its future core descendants, including itself, and
put
\[
 G(t)=\sup_xF_x(t),\qquad H(t)=1+\int_0^tG(u)\,du.
\]
The admissible class and its harmless finite inert-boundary completion
are specified in Section~\ref{model:section}. In the inert newborn
process, let $\sigma_v$ be the birth time of core $v$, and give the
initial core genealogical depth zero. Define
\[
 A_m^\infty(s)=\mathbb E_\infty
   \sum_{\operatorname{depth}(v)=m}e^{-s\sigma_v},\qquad s>0.
\]

\begin{theorem}[Exact characterization and effective computation]
\label{main:computation}
The limits
\[
 r_\infty(s)=\lim_{m\to\infty}A_m^\infty(s)^{1/m}
 \quad(s>0)
\]
exist and define a positive, continuous, strictly decreasing function.
Its unique crossing of one is the original friend-tree exponent:
\[
 r_\infty(\mu)=1.
\]
Moreover,
\[
 \lim_{t\to\infty}\frac{\log G(t)}t
 =\mu
 =\inf_{T\in\mathbb Q_{>0}}\frac{\log H(T)}T.
\]
Every fixed coefficient $A_m^\infty(s)$ at rational $s>0$, and every
$H(T)$ at rational $T>0$, is effectively computable. There is an
algorithm which, on input a positive rational $\varepsilon$, halts
with rational $a<b$ satisfying
\[
 a<\mu<b,\qquad b-a<\varepsilon.
\]
\end{theorem}

The proof combines Theorems~\ref{gc:root}, \ref{calc:calculator},
\ref{erc:equality}, and \ref{ecc:computable}. The lower certificate is
particularly explicit:
\[
 s<\mu\quad\Longleftrightarrow\quad
 A_m^\infty(s)>12(s+3)\quad\hbox{for some finite }m.
\]
Finite positive rational-history sums eventually witness every strict
lower inequality. The submultiplicative time-block formula supplies
the complementary complete search for strict upper inequalities.
The resulting algorithm has a proof of termination and a prescribed
error, but no useful complexity bound is established here.

For a retained state $y$, let $P_t^\infty$ be the inert semigroup and
write
\[
 V_s(y)=\int_0^\infty e^{-st}P_t^\infty C(y)\,dt,
 \qquad
 W(s)=\sup_x\int_0^\infty e^{-st}F_x(t)\,dt.
\]
If $y^+$ is obtained by adding a leaf at the retained root, the main
uniform estimate is as follows.

\begin{theorem}[Boundary gradient and robustness]
\label{main:boundary}
For every $s>\mu$, the preceding quantities are finite, and
\[
 |V_s(y^+)-V_s(y)|\le\frac{40}{s}W(s)
\]
uniformly over all finite retained states $y$. A finite time-block
certificate computes an explicit upper bound for the right side.
Fix $Q<\infty$ and permit an arbitrary nonanticipating root-leaf input
with predictable rate $0\le q_t\le Q$. Starting from the inert newborn,
let $F_Q^-(t)$ and $F_Q^+(t)$ be the infimum and supremum of the expected
core counts over these policies. Then
\[
 \lim_{t\to\infty}\frac{\log F_Q^-(t)}t
 =\lim_{t\to\infty}\frac{\log F_Q^+(t)}t=\mu.
\]
The conclusion includes policies selected separately for each horizon,
provided that the same bound $Q$ applies.
\end{theorem}

The finite-horizon gradient is proved before uniform resolvent
finiteness is known. Its proof uses a disagreement exploration with
discounted reproduction mass at most $19/20$ at $s=11/20$. This avoids
assuming the conclusion in the comparison. Theorems~\ref{erc:gradient}
and~\ref{bcl:matching} give the precise statements and proofs.

We also refine the existing rational motif certificate.

\begin{theorem}[A certified numerical refinement]
\label{main:bound}
The original exponent satisfies
\[
 \frac{113}{200}\le\mu\le
 \frac{2832242711}{5000000000}=0.5664485422.
\]
The lower endpoint is imported from~\cite{Uzay}. The upper endpoint
follows from an exact rational modification of its upper motif
potential, with a positive reserve in every vertex row.
\end{theorem}

This is a refinement of an existing computer-assisted certificate.
It reduces the interval width by approximately $3.43\%$; it does not
give high-precision digits of $\mu$. Section~\ref{numerical:section}
states the all-degree extension and the exact verification procedure.

\paragraph{Scope of the results.}
This paper proves mean exponential growth statements, exact
characterizations, and computability, using the stated exponent
existence input. It does not prove a sharp asymptotic
$\mathbb E C_n\sim c n^\mu$, a positive almost-sure amplitude, or
irrationality or transcendence of $\mu$. A separate companion draft
studies independent branching representations and irrationality for
modified local redirection models~\cite{Local}; those exponents are
distinct from the original exponent studied here.

\paragraph{Organization.}
Sections~\ref{model:section}--\ref{calc:section} give the exact kernel,
inert comparison, pressure equation, and finite coefficient calculator.
The uniform mean and locality estimates are followed by finite-window
supersolutions, controlled moments, and the bounded-parent likelihood
argument. Section~\ref{erc:section} proves the boundary gradient;
Section~\ref{erc:closure-section} identifies the uniform exponent.
The final sections derive computability, construct the resolvent,
prove boundary robustness, and give the numerical certificate.

\section{The core process and exact waiting-event elimination}
\label{model:section}

\subsection{States, genealogy, and the imported exponent}

Each real vertex has an independent rate-one source clock. At a ring it
chooses a neighbour uniformly and creates a new leaf at that neighbour.
The total birth rate of a full tree of size $n$ is $n$, so its embedded
arrival chain is the discrete random friend tree. Every finite tree is
reachable from an initial edge with positive probability in a fixed
finite time. Consequently the expected core count from any fixed
finite seed has upper exponential rate at most the rate from the edge.

We use the following established input from~\cite{Uzay}: for the
process started from an edge,
\begin{equation}\label{model:imported-exponent}
 F_2(t):=\mathbb E_{\mathrm{edge}}C(t)=e^{\mu t+o(t)},
 \qquad
 \mathbb E C_n=n^{\mu+o(1)},
 \qquad 0.5650\le\mu\le0.5665.
\end{equation}
The proof source is fixed at the revision in the bibliography. The
weighted planted-branch comparison used below is also stated explicitly
with its location in that source. We do not use an unweighted ordering
between planted and free branch means.

A core state consists of a finite tree $H$ of real cores, leaf counts
$L_i\ge0$, and inert-stub counts $b_i^\partial\ge0$. Put
\[
 k=|H|,\qquad c_i=\deg_H(i)+b_i^\partial,
 \qquad D_i=L_i+c_i\ge2.
\]
An inert stub contributes to $D_i$ but has no clock. A choice of that
stub produces a null transition. For a full tree all stub counts are
zero. In that case a singleton core has at least two leaves. In the
one-stub retained newborn, $k=1$, $c_1=L_1=1$, so the singleton has
exactly one real leaf. These two cases have different core-ring
transition laws and are kept separate.

Whenever a core chooses one of its leaves, that leaf becomes a new
core with one leaf of its own. Call it a core child of the selecting
core. Other arrivals only add a leaf to an existing core. This defines
a permanent genealogy of cores. A newborn tag has one leaf and one
core neighbour, either a real parent or an inert stub. Let $\mathcal X$
be the countable class of all finite environments with such a tagged
newborn, allowing finitely many inert stubs away from the tag as well.
Only the tag and its future genealogical descendants are counted in
$F_x(t)$.

For $s>0$, define a positive operator on bounded functions on
$\mathcal X$ by
\begin{equation}\label{model:tagged-kernel}
 (\mathcal K_s f)(x)
 =\mathbb E_x\sum_{v:\,v\text{ a future child of the tag}}
             e^{-s\sigma_v} f(X_v),
\end{equation}
where $X_v$ is the entire environment at the birth of $v$, retagged
at $v$. The strong Markov property and Tonelli's theorem give
\[
 (\mathcal K_s^m\mathbf1)(x)
 =\mathbb E_x\sum_{\operatorname{depth}(v)=m}e^{-s\sigma_v}
 =A_m(s;x).
\]
The whole environment is retained in $X_v$; this is not an assertion
that descendant families evolve independently.

Give every potential core a rate-one clock and let every ring create a
potential child, whether or not the corresponding actual ring promotes
a leaf. Assign unused potential branches fresh independent drivers.
Every actual tagged genealogy is a subset of this potential genealogy,
whose expected depth-$j$ count at time $t$ is $t^j/j!$. In particular
\begin{equation}\label{model:kernel-bounds}
 \|\mathcal K_s\|_\infty\le1/s,
 \qquad [2(s+2)]^{-m}\le A_m(s;x)\le s^{-m}.
\end{equation}
For the lower bound, wait for the first ring of the tag or its initial
leaf, and require a tag ring choosing a leaf. Before that ring the
tag still has one core neighbour, while outside arrivals can only add
leaves. Its leaf-choice probability is at least $1/2$. The discounted
probability is therefore at least $1/[2(s+2)]$. Iterate at newborn
states. Core-count Yule domination also gives $F_x(t)\le e^t$ uniformly.

Finite inert environments introduce no larger fixed-time supremum.
Replace each of their $b$ stubs by an edge to a genuine auxiliary
parent of degree $R$, completed by $R-1$ leaves. Until a return across
one of these edges, the retained histories agree. The probability of
such a return by a fixed time $T$ is at most $bT/R$, since degrees
never decrease. Uniform tagged Yule second moments and
Cauchy--Schwarz make the difference in tagged expected rewards tend
to zero as $R\to\infty$. Thus every fixed inert profile is a
finite-horizon limit of genuine finite environments.

\subsection{The integrated leaf burst}

Between two core-clock rings, the core skeleton is fixed. Each leaf
creates leaves at its own core at rate one, so the coordinates $L_i$
are independent Yule processes. The next core ring has an independent
exponential waiting time of rate $k$. Write $\mathcal Y_t$ for the
leaf semigroup and set
\begin{equation}\label{cc:R-def}
 (R_{H,s}f)(L)=\int_0^\infty e^{-(k+s)t}(\mathcal Y_t f)(L)\,dt,
 \qquad \alpha=k+s,\quad \ell=\sum_iL_i.
\end{equation}
If $I=\{i:L_i>0\}$ and $\ell>0$, the exact matrix entry is
\begin{equation}\label{cc:beta-entry}
 R_{H,s}(L,M)=
 \left[\prod_{i\in I}\binom{M_i-1}{L_i-1}\right]
 \mathrm B(\alpha+\ell,|M|-\ell+1)
\end{equation}
when $M_i\ge L_i$ for $i\in I$ and $M_i=0$ otherwise, and it is zero
off this support. If $\ell=0$, the leaf vector remains zero and the
only entry is $R_{H,s}(0,0)=1/\alpha$.

Indeed, with $z=e^{-t}$, the Yule transition probability from $L_i>0$
to $M_i$ is
$\binom{M_i-1}{L_i-1}z^{L_i}(1-z)^{M_i-L_i}$.
Multiplication and integration against $dt=-dz/z$ give
\eqref{cc:beta-entry}. Thus leaf-clock waiting events are integrated
exactly, including their unbounded population growth.

At rational $s$, every entry is rational, since
\[
 \mathrm B(a,b+1)=\frac{b!}{a(a+1)\cdots(a+b)}
 \qquad(a>0,\ b\in\mathbb Z_{\ge0}).
\]
The exact total-leaf tail, for $N\ge\ell>0$, is
\begin{equation}\label{cc:exact-tail}
 \alpha\sum_{|M|\ge N}R_{H,s}(L,M)
 =\prod_{j=\ell}^{N-1}\frac{j}{j+\alpha}.
\end{equation}
To see this, observe the total-leaf Yule process at an independent
rate-$\alpha$ exponential time. Its hitting time of $N$ is the sum
of independent exponentials of rates $\ell,\ldots,N-1$.

\subsection{The exact core-ring kernel}

At a core ring with source $j$, every real core neighbour $i$ receives
a leaf with probability $1/(M_j+c_j)$; an inert choice is null; and
a promotion occurs with probability $M_j/(M_j+c_j)$. Define the rate
operators
\begin{align}
 (A_Hf)(M)
 &=\sum_j\left[
   \sum_{i\sim_Hj}\frac{f(M+e_i)}{M_j+c_j}
   +\frac{b_j^\partial f(M)}{M_j+c_j}\right],
   \label{model:A}\\
 (B_Hg)(M)
 &=\sum_{j:M_j>0}\frac{M_j}{M_j+c_j}
      g(H^{(j)},M^{(j)}).
   \label{model:B}
\end{align}
Here $H^{(j)}$ adjoins a child at $j$; $M^{(j)}$ decreases $M_j$ by one
and appends the new child's leaf count one. The stub counts are inherited.
Their row masses sum to $k$. Hence
\begin{equation}\label{model:Q}
 Q_s=R_s(A+B),\qquad Q_s\mathbf1(y)=\frac{C(y)}{C(y)+s}
\end{equation}
is the exact discounted transition kernel between consecutive core
rings. The notation $R_s$ denotes the appropriate $R_{H,s}$ in each
state. Null stub choices are included and consume one ring.

If desired, nonpromotion core rings can also be eliminated. With
$W_{H,s}=R_{H,s}A_H$ and $V_{H,s}=R_{H,s}B_H$, the exact first-promotion
kernel is
\begin{equation}\label{model:first-promotion}
 \mathcal J_{H,s}=\sum_{q\ge0}W_{H,s}^qV_{H,s}.
\end{equation}
Before the first promotion, the promotion intensity
$b_H(M)=\sum_jM_j/(M_j+c_j)$ does not decrease. Thus on the upward
cone above $L$,
\[
 W_{H,s}\mathbf1\le\frac{k-b_H(L)}{k+s}<1\qquad(s>0).
\]
The omitted mass after $q=0,\ldots,Q-1$ is at most
$[(k-b_H(L))/(k+s)]^Q$. The separate total-leaf tail
\eqref{cc:exact-tail} makes every finite calculation effective.
We use the full ring kernel $Q_s$ when genealogical marks are required;
its histories have explicit nonnegative rational weights at rational
$s$. Section~\ref{calc:section} gives a complete coefficient calculator.

\section{Uniform comparison with an inert parent}

Let $\mathcal X$ be the class of finite newborn birth environments,
allowing finitely many inert core-neighbour stubs. The tag $u$ has one
leaf and one core neighbour, which is either a real parent or an inert
stub. Write
\[
 A_m(s;x)=(\mathcal K_s^m\mathbf1)(x),\qquad
 A_m^\infty(s)=A_m(s;x_\infty),
\]
where $x_\infty$ is the tag with its leaf and an inert parent stub.
All generations count only future core descendants of the tag.
The potential-genealogy construction in Section~\ref{model:section} gives
$\|\mathcal K_s\|_\infty\le1/s$ and
$A_m(s;x)\ge[2(s+2)]^{-m}$.

\begin{theorem}[Uniform boundary removal]\label{growth:minorization}
For every $s>0$, every $m\ge1$, and every $x\in\mathcal X$,
\begin{equation}\label{growth:comparison}
 \boxed{
 A_m(s;x)\ge\gamma_s A_m^\infty(s),
 \qquad \gamma_s=\frac1{12(s+3)}.
 }
\end{equation}
The constant is independent of both the initial environment and the
genealogical generation.
\end{theorem}

The point of the proof is to create a permanent no-return event using
only exterior randomness. Merely conditioning on the absence of actual
parent arrivals would generally bias the descendant process.

\subsection{An exterior process with a positive permanent success probability}

Consider a birth--death process $X$ on the nonnegative integers,
absorbed at zero, with births $x\to x+1$ at rate $x$ and deaths at
rate one while $x>0$. At every death-clock ring, attach an independent
uniform mark $U$. Kill the construction if $U\le1/X(t-)$; otherwise
carry out the death. A ring at $X=1$ always kills it. The generator of
the surviving process, acting on functions with value zero at the
cemetery state, is
\begin{equation}\label{growth:killed-generator}
 Gh(x)=x\bigl(h(x+1)-h(x)\bigr)
 +(1-1/x)\bigl(h(x-1)-h(x)\bigr)-h(x)/x.
\end{equation}

\begin{lemma}\label{growth:survival}
Starting from $X(0)=1$, the probability of never being killed is at least
$1/6$.
\end{lemma}
\begin{proof}
Use the bounded nonnegative function
\[
 h(0)=0,\qquad h(1)=\frac16,\qquad
 h(x)=1-\frac4{3x}\quad(x\ge2).
\]
Direct substitution gives
\[
 Gh(1)=0,\qquad Gh(2)=\frac7{36},\qquad
 Gh(x)=\frac{x-3}{3x(x+1)}\ge0\quad(x\ge3).
\]
The process is nonexplosive, being dominated by a linear pure-birth
process. Dynkin's formula, or localization followed by bounded
convergence, gives
$\mathbb E[h(X(t));\,\text{not killed by }t]\ge h(1)=1/6$.
Since $0\le h\le1$, survival to every finite $t$ has probability at least
$1/6$. Let $t\to\infty$.
\end{proof}

\subsection{Proof of the uniform comparison}

If the parent is already inert, the tagged component has exactly the
inert-parent law, so~\eqref{growth:comparison} is immediate.
Suppose the parent $p$ is real.

First prepare a parent leaf without changing the tag. Wait for the first
ring among the tag, its sole leaf, and its parent. These are three
independent rate-one clocks. Declare preparation successful only when
the first ring is the tag's ring and it chooses the parent, rather than
the leaf. Until that ring the tag remains in its newborn state, even
if the parent receives arrivals from elsewhere. If $T$ is this first
ring time and $E_0$ the successful outcome, then
\begin{equation}\label{growth:preparation}
 \mathbb E[e^{-sT}\mathbf1_{E_0}]
 =\int_0^\infty \tfrac12 e^{-(s+3)t}\,dt
 =\frac1{2(s+3)}.
\end{equation}
On $E_0$ no tagged child has yet been born, the tag still has its one
leaf and one parent, and the parent has gained a leaf.

Restart the clocks at $T$. Lower-couple the parent's leaf count $L_p$
with the exterior process $X$ of Lemma~\ref{growth:survival}, starting
from one. Use a common parent-leaf birth Poisson random measure:
marks below $X$ create births in $X$, and marks below $L_p$ create the
actual parent-leaf births. This measure is independent of all future
tag-descendant driving randomness. At \emph{every} parent core-clock
ring, decrease $X$ by one while it is positive; the actual $L_p$
decreases by at most one. All other arrivals at the parent only increase
$L_p$. Therefore, up to killing,
\[
 L_p(t)\ge X(t),\qquad D_p(t)\ge X(t).
\]

At a parent core-clock ring, choose the tag as target precisely when
$U\le1/D_p(t-)$. Otherwise choose uniformly among the remaining current
neighbours. This is an exact uniform-neighbour construction. The
reserved event $U\le1/X(t-)$ contains every possible parent-to-tag
arrival. The process $X$ uses the parent ring times but not these marks.

Let $E$ be the event that the exterior construction is never killed.
It depends only on the parent-leaf birth measure, the parent core-clock
times, and the reserved marks. In particular, conditional on the
preparation history, it is independent of the future driving randomness
inside the tag's descendant component. By the lemma, its conditional
probability is at least $1/6$.

On $E$, the parent never sends a leaf to the tag. The tag may still
send leaves to the parent, but such outward events have no effect on
its retained component and cannot create a return arrival on $E$.
Its entire descendant history therefore has the inert-parent law,
using the unchanged independent interior drivers. Positivity, the
strong Markov property at $T$, and~\eqref{growth:preparation} give
\[
 A_m(s;x)
 \ge
 \mathbb E[e^{-sT}\mathbf1_{E_0}]\,
 \frac16\,A_m^\infty(s)
 =
 \frac{A_m^\infty(s)}{12(s+3)}.
\]
This proves Theorem~\ref{growth:minorization}. The same argument applies
when other exterior neighbours are inert stubs: it only uses a real
parent's clock, its own leaf births, and nonnegative additional
immigration.

\section{An actual generation limit for the inert-parent process}

Define the uniform lower and upper sequences
\[
 u_m(s)=\inf_{x\in\mathcal X}A_m(s;x),\qquad
 v_m(s)=\sup_{x\in\mathcal X}A_m(s;x).
\]
The inert-parent state belongs to $\mathcal X$, so
\begin{equation}\label{growth:inf-comparison}
 \gamma_s A_m^\infty(s)\le u_m(s)\le A_m^\infty(s).
\end{equation}

\begin{theorem}\label{growth:inert-limit}
The limit
\[
 r_\infty(s)=\lim_{m\to\infty}A_m^\infty(s)^{1/m}
\]
exists, is finite and positive, and equals
$\lim_{m\to\infty}u_m(s)^{1/m}$. More precisely,
\begin{equation}\label{growth:sup-formula}
 \boxed{
 r_\infty(s)
 =\sup_{m\ge1}u_m(s)^{1/m}
 =\sup_{m\ge1}\bigl(\gamma_s A_m^\infty(s)\bigr)^{1/m}.
 }
\end{equation}
For every $m\ge1$,
\begin{equation}\label{growth:uniform-inert-upper}
 A_m^\infty(s)\le\gamma_s^{-1}r_\infty(s)^m.
\end{equation}
\end{theorem}
\begin{proof}
Positivity of the kernel gives
\[
 u_{m+n}\ge u_m u_n.
\]
Also, applying~\eqref{growth:comparison} to the output birth environments
of the first $m$ tagged transfers gives
\[
 A_{m+n}^\infty
 =(\mathcal K_s^m A_n)(x_\infty)
 \ge\gamma_s A_m^\infty A_n^\infty.
\]
Thus both $u_m$ and $\gamma_s A_m^\infty$ are supermultiplicative.
Their logarithms have finite linear bounds, because
$[2(s+2)]^{-m}\le A_m(s;x)\le s^{-m}$.
Fekete's lemma and~\eqref{growth:inf-comparison} prove the common limit
and both supremum formulae. The second supremum formula immediately
gives~\eqref{growth:uniform-inert-upper}.
\end{proof}

The theorem supplies a complete scalar certificate criterion for
$r_\infty(s)>1$:
\[
 r_\infty(s)>1
 \quad\Longleftrightarrow\quad
 \gamma_s A_m^\infty(s)>1
 \quad\text{for some finite }m.
\]
At rational $s$, $\gamma_s$ is rational and the existing finite-history
approximations of $A_m^\infty(s)$ can be used in exact lower certificates.
The next section identifies the inert continuous-time growth with the original exponent.


\section{The inert process has the original mean exponent}
\label{gs:inert-exponent-section}

Let $\mathcal X$ be the class of finite birth environments, allowing
inert boundary stubs, with the tag initially satisfying $L_u=c_u=1$.
For $x\in\mathcal X$, let $Z_m(t;x)$ count its future core descendants
of genealogical generation $m$ born by time $t$, and put
\[
 F_x(t)=1+\sum_{m\ge1}\mathbb E_x Z_m(t;x).
\]
The initial tag contributes the term $1$; other initial core vertices
are not counted. Write $x_\infty$ for the tag with one leaf and one
inert parent stub, and $F_\infty=F_{x_\infty}$.

The input from~\cite{Uzay} is the original free-tree mean exponent
\[
 F_2(t):=\mathbb E_{\mathrm{edge}} C(t)
       =e^{\mu t+o(t)}.
\]
We also use the planted-branch comparison proved in its
\texttt{proofs/existence.md} in Lemma 4.1(c), Corollary 5.3, and Lemma 6.1 of that source. We state precisely the portion used here.
Fix $M\ge2$, sufficiently small $\epsilon>0$, and
$k\ge6M/\epsilon$. Put
\[
 \eta=b(\epsilon)\log2+\frac9M,
 \qquad b(\epsilon)=7/\log(1/\epsilon)\longrightarrow0
       \quad(\epsilon\downarrow0).
\]
For all sufficiently large $t$, there is a rooted $k$-vertex free tree
$(T,v)$ for which the following lower comparison holds. If $T$ is
planted as a branch with its root $v$ attached to an exterior vertex
of full degree at least $M$, its future expected internal core count
is at least
\begin{equation}\label{gs:planting-input}
 \frac1k e^{-\eta t}
       \bigl[F_2(t)-k-O(e^{-t})\bigr].
\end{equation}
The exterior may evolve freely. The constants are uniform over the
finite collection of rooted $k$-vertex seeds. The seed and root can
depend on $t$, but are selected from their laws before planting.

For clarity, the comparison weights a free reference history by
\[
 e^{-t/M}\prod_{\tau\le t:\,\text{source }v}
       \frac{d_v(\tau-)}{d_v(\tau-)+1},
\]
where $d_v$ is its internal reference degree. A high-degree original
vertex in the reference history gives the factor $e^{-\eta t}$;
averaging over its $k$ possible choices gives $1/k$. The inequality
$F_2(t)\le k+\max_{|T|=k}\mathbb E C_T(t)$ supplies the displayed
free-tree term. This is the cited comparison input, not a new
independence assumption about the two sides of an edge.

\begin{lemma}[Uniform finite preparation]\label{gs:uniform-preparation}
For fixed $M,k$, there are $h<\infty$ and $p>0$ such that, from every
$x\in\mathcal X$, any chosen rooted $k$-vertex seed $(T,v)$ can be
planted by time $h$ with probability at least $p$. Every core vertex in
the branch belongs to the original tag's future core genealogy, and its
root is attached to that tag, whose full degree is then at least $M$.
Both constants are uniform
in the exterior environment and in the rooted seed.
\end{lemma}

\begin{proof}
If the tag has a real parent, reserve potential parent-to-tag arrows
at rate $1/2$ and accept each with probability $2/D_p(t)$. The parent's
full degree is at least two and never decreases, so this is an exact
construction. Suppressing these potential arrows during a fixed
preparation interval costs only $e^{-h/2}$, independently of the
retained component's clocks. For a stub parent the same suppression
can be imposed on an independent unused stream. The preparation can
therefore be performed in the inert-parent process.

First keep the tag silent for one unit of time and let its one leaf
produce enough Yule offspring that the tag's full degree is at least
$M$. This has positive probability. In a second unit interval require
exactly one tag ring, choosing a leaf, and keep the resulting two-vertex
branch silent until that interval ends. The leaf-choice probability is
at least $(M-1)/M$, since all growth before that promotion increases
the number of leaves at the tag. Its full degree is unchanged by the
promotion.

Starting with that branch edge, build the specified rooted tree by
adjoining its remaining vertices in an order in which each new vertex
has an already present parent. During each unit interval keep the tag
silent and prescribe one ring of an internal neighbour of the desired
attachment vertex, its choice, and silence of all other internal clocks.
Also require the newborn to remain silent until the interval ends.
All relevant internal degrees are at most $k+1$, so every prescribed
choice has probability at least $1/(k+1)$. Old leaves outside the branch
can grow freely; they only increase the exterior tag's degree.

There are finitely many prescribed clocks and choices and finitely many
rooted seeds of size $k$. Taking their minimum probability gives
$p>0$, with a common preparation time, for example $h=k$.
\end{proof}

\begin{theorem}[Uniform lower growth and the inert exponent]
\label{gs:inert-mu}
For every $\varepsilon>0$ there are $c_\varepsilon>0$ and
$t_\varepsilon<\infty$ such that
\begin{equation}\label{gs:uniform-lower-time}
 \inf_{x\in\mathcal X}F_x(t)
       \ge c_\varepsilon e^{(\mu-\varepsilon)t}
       \qquad(t\ge t_\varepsilon).
\end{equation}
Moreover the inert process has the exact mean exponent
\begin{equation}\label{gs:inert-mean-limit}
 \lim_{t\to\infty}\frac{\log F_\infty(t)}t=\mu.
\end{equation}
\end{theorem}

\begin{proof}
After the uniform preparation, apply~\eqref{gs:planting-input} to its
fresh branch. Every reference core counted by that comparison is an
actual future descendant of the original tag. Thus, for all large $t$,
\begin{equation}\label{gs:uniform-planted-bound}
 \inf_{x\in\mathcal X}F_x(h+t)
 \ge\frac pk e^{-\eta t}
          \bigl[F_2(t)-k-O(e^{-t})\bigr].
\end{equation}
Uniformity of $p,h$ over the finite seed set permits the optimizing
seed to vary with $t$. Fix the comparison parameters before taking
$t\to\infty$, and then let $\eta\downarrow0$. This proves
\eqref{gs:uniform-lower-time}, and in particular the required lower
bound for the exponent of $F_\infty$.

For the upper bound, embed the inert tag and its leaf in a finite real
tree $T_M$ whose parent has full degree $M$: give the parent $M-1$
additional leaves. Reserve its potential arrow to the tag at rate
$1/M$, accepting with probability $M/D_p(t)$. Before the first
potential arrow, the retained tag component follows exactly the inert
law; a tag ring choosing the exterior affects only the parent and
does not change the retained component. The potential-arrow stream is
independent of the inert component's drivers. Therefore
\begin{equation}\label{gs:inert-upper-embedding}
 e^{-t/M}F_\infty(t)\le\mathbb E C_{T_M}(t).
\end{equation}
Every fixed finite tree is reachable from the original edge by a fixed
time with positive probability. The Markov property consequently gives
constants $a_M,b_M<\infty$ with
\[
 \mathbb E C_{T_M}(t)\le a_MF_2(t+b_M).
\]
For fixed $M$,~\eqref{gs:inert-upper-embedding} implies
$\limsup t^{-1}\log F_\infty(t)\le\mu+1/M$. Letting $M\to\infty$
finishes the proof.
\end{proof}


\section{A scalar growth limit with the exact critical root}
\label{gc:section}

The inert comparison gives an actual generation limit, and the preceding continuous-time comparison identifies its critical parameter.

\begin{theorem}[Exact inert-parent characteristic equation]
\label{gc:root}
For $s>0$, let
\[
 r_\infty(s)=\lim_{m\to\infty} A_m^\infty(s)^{1/m}.
\]
Then $r_\infty$ is positive, continuous, and strictly decreasing, and
\begin{equation}\label{gc:crossing}
 \boxed{r_\infty(\mu)=1,\qquad
 r_\infty(s)>1\ (0<s<\mu),\qquad
 r_\infty(s)<1\ (s>\mu).}
\end{equation}
In particular the original exponent is the unique solution of this
scalar equation. For every $m\ge1$,
\begin{equation}\label{gc:critical-bound}
 A_m^\infty(\mu)\le 12(\mu+3)<43.
\end{equation}
\end{theorem}
\begin{proof}
Tonelli's theorem gives the extended-value identity
\[
 1+\sum_{m\ge1}A_m^\infty(s)
 =s\int_0^\infty e^{-st}F_\infty(t)\,dt.
\]
By Theorem~\ref{gs:inert-mu}, this is finite for $s>\mu$ and infinite
for $s<\mu$. Hence $r_\infty(s)\le1$ above $\mu$ and
$r_\infty(s)\ge1$ below $\mu$.

Each $A_m^\infty$ is a Laplace transform of a positive measure, so
$m^{-1}\log A_m^\infty(s)$ is convex. Its pointwise limit
$g(s)=\log r_\infty(s)$ exists by Theorem~\ref{growth:inert-limit}
and has the finite bounds
\[
 -\log\bigl(2(s+2)\bigr)\le g(s)\le-\log s.
\]
Thus $g$ is finite, convex, continuous, and nonincreasing on
$(0,\infty)$. If it took the same value at two distinct points,
convexity and monotonicity would force it to be constant thereafter.
This contradicts $g(s)\le-\log s\to-\infty$. It is strictly decreasing.
The weak inequalities on either side of $\mu$, together with continuity
and strict decrease, prove~\eqref{gc:crossing}.
Equation~\eqref{growth:uniform-inert-upper} at $s=\mu$ gives
\eqref{gc:critical-bound}; the numerical inequality uses the imported
upper bound $\mu\le1133/2000$.
\end{proof}

\begin{corollary}[Complete scalar lower certificates]
\label{gc:scalar-certificates}
For every $s>0$,
\begin{equation}\label{gc:certificate-equivalence}
 \boxed{s<\mu\quad\Longleftrightarrow\quad
 A_m^\infty(s)>12(s+3)\text{ for some finite }m.}
\end{equation}
At rational $s$, finite positive rational-history approximations provide
a complete semidecision procedure for $s<\mu$. The same statement
holds for a positive algebraic $s$ specified exactly: each finite
comparison is then an exact algebraic-number sign calculation.
\end{corollary}
\begin{proof}
Combine~\eqref{growth:sup-formula} with~\eqref{gc:crossing}.
When the inequality holds, its strict margin is eventually retained
by finite positive history approximations. Conversely every successful
finite lower approximation is bounded by the exact quantity and is
therefore sound. All its weights are rational functions of $s$ with
positive denominators for $s>0$.
\end{proof}

For all sufficiently large $m$, the equation
\[
 \frac{A_m^\infty(s_m)}{12(s_m+3)}=1
\]
has a unique positive solution and $s_m\le\mu$, with $s_m\to\mu$.
Indeed, for each fixed $s<\mu$, the left side at $s$ eventually exceeds
one by the actual generation limit; at $\mu$ it is at most one and at
infinity it tends to zero. Each individual left side is continuous and
strictly decreasing. The running maxima of these roots therefore
increase to $\mu$. No effective outer error or monotonicity of the
individual $s_m$ is asserted.


\section{A finite rational calculator for the inert coefficients}
\label{calc:section}

This section gives a direct computation of $A_m^\infty(s)$, without
enumerating all pruned environment profiles. The two budgets count
different events: $R$ bounds \emph{all} actual core-clock rings, including
null choices of the inert parent, while $N$ bounds real vertex arrivals
after the initial tag and its one leaf. Throughout, $m\ge1$ and $s>0$.

\begin{lemma}[A universal core-ring tail]\label{calc:ring-tail}
Put
\[
 P_R(z)=\prod_{i=1}^R\left(1+\frac zi\right),
 \qquad P_0(z)=1.
\]
The expected discounted generation-$m$ births occurring after the $R$th
actual core-clock ring in the inert process are at most
\begin{equation}\label{calc:tail-formula}
 T_{m,R}(s)=
 \left(\prod_{i=1}^R\frac{i}{i+s}\right)
 \sum_{j=0}^{m-1}[z^j]P_R(z)\,s^{-(m-j)}.
\end{equation}
This is rational for rational $s$. For fixed $m,s$ it tends to zero,
with $T_{m,R}(s)=O_{m,s}(R^{-s}(1+\log R)^{m-1})$.
\end{lemma}

\begin{proof}
Use the common potential genealogy in which every ring of every
potential core creates a potential child, even if its actual counterpart
has a null or nonpromotion event. Actual cores and their clock rings
are subsets of this potential process. The potential core population
is a rate-one Yule process started with one vertex. Write $\tau_R$ for
its $R$th birth time. The $R$th actual core ring occurs no earlier than
$\tau_R$, and every omitted actual generation-$m$ birth is a potential
generation-$m$ birth after $\tau_R$.

At potential size $i$, the holding time has rate $i$, independently of
the uniform choice of its parent. Thus the generation profile at
$\tau_R$ is independent of $\tau_R$, and
\[
 \mathbb E e^{-s\tau_R}=\prod_{i=1}^R\frac{i}{i+s}.
\]
Its expected profile polynomial is $P_R$: adding a vertex at size $i$
multiplies the previous expected polynomial by $1+z/i$.

Conditional on the profile at $\tau_R$, a currently existing vertex
at depth $j<m$ contributes at most, and in the potential process exactly,
$s^{-(m-j)}$ from its new future branches. These contributions do not
double count existing child branches: a future vertex is assigned to
its deepest ancestor already present at $\tau_R$. Summing, and using
independence, proves~\eqref{calc:tail-formula}.
Finally $[z^j]P_R\le H_R^j/j!$, where $H_R=\sum_{i=1}^R1/i$, and the
product in~\eqref{calc:tail-formula} is $O_s(R^{-s})$.
\end{proof}

\subsection{Bounding the histories discarded by a birth budget}

The retained state consists of a finite rooted core, its leaf counts,
and the initial inert stub. Its core depths are the genealogical
depths from the original tag. For such a state $x$, define
\begin{equation}\label{calc:future-reward}
 V_m(x)=\sum_{v:\,d(v)<m}s^{-(m-d(v))},
 \qquad C_m(s)=\max_{0\le r<m}s^{-r}.
\end{equation}
The potential-genealogy comparison shows that $V_m(x)$ bounds all
future generation-$m$ reward from that state, independent of its
leaf counts.

Between core rings, integrate the leaf growth using the exact beta
burst kernel. If a burst is discarded because it exceeds the real-birth
budget, no core birth occurs before its next forced core ring, and its
core graph and depths are unchanged. That ring can contribute at most
one new core to the future reward, or one immediate generation-$m$
birth. Consequently the entire omitted reward from this branch is at
most its discounted omitted-ring mass times
\begin{equation}\label{calc:omitted-burst}
 V_m(x)+C_m(s).
\end{equation}
The omitted-ring mass is computed exactly from the Yule burst tail,
including its factor $k/(k+s)$ for $k$ current cores.

For a retained burst whose subsequent real core event itself exceeds
the birth budget, use the sharper bound
\begin{equation}\label{calc:omitted-core-event}
 \mathbf1_{\{\text{new core has depth }m\}}+V_m(x_{\mathrm{post}}),
\end{equation}
weighted by that event's exact discounted mass. A nonpromotion event
has no immediate reward. Null stub choices use one core ring but no
real birth and are retained accordingly.

\begin{proposition}[Certified scalar intervals]\label{calc:calculator}
Enumerating beta bursts through at most $R$ core rings and at most $N$
real births gives a positive finite rational sum $L_{m,R,N}(s)$ at
rational $s$. Let $E_{m,R,N}(s)$ be the sum of the discarded-history
bounds~\eqref{calc:omitted-burst}--\eqref{calc:omitted-core-event}. Then
\begin{equation}\label{calc:interval}
 L_{m,R,N}(s)\le A_m^\infty(s)
 \le L_{m,R,N}(s)+E_{m,R,N}(s)+T_{m,R}(s).
\end{equation}
The upper endpoint may also be capped at $s^{-m}$. These bounds compute
$A_m^\infty(s)$ to any prescribed rational accuracy: first increase
$R$ until its explicit tail is small, then increase $N$ until the
computed discarded-history bound is small.
\end{proposition}

\begin{proof}
Retained histories have their exact positive discounted weights and
are a subset of the original histories. The first discarded burst or
core event accounts for a missing path by
\eqref{calc:omitted-burst} or~\eqref{calc:omitted-core-event}. The additional
universal tail bounds all births after $R$ actual core rings; any
overlap between the two error terms is harmless. This proves
\eqref{calc:interval}.

For fixed $R$, at most $R+1$ real cores exist through these rings, so
the reward bounds at every first discarded history are bounded by a
constant depending only on $R,m,s$. The first-discard histories are
disjoint. Their total undiscounted probability is at most the
probability that the real population at the $R$th core ring exceeds
$N+2$. That population is finite almost surely: the $R$th core ring
occurs by the $R$th ring of the original tag, and the vertex process
is nonexplosive, with birth rate at most its current population.
Thus $E_{m,R,N}(s)\to0$ as $N\to\infty$ for fixed $R$.

Every step is finite: before each retained ring there are only
finitely many leaf allocations within the birth budget, and null
events consume the ring budget. At rational $s$, each beta entry,
neighbour-choice rate, tail, and reward bound is rational.
\end{proof}

\subsection{Implementation and the complete scalar lower test}

The standard-library program \texttt{inert\_growth\_kernel.py} imports
the exact burst routines from \texttt{core\_clock\_kernel.py}. It keeps
the original tag at root zero and combines root-isomorphic states.
It reports the lower and upper coefficient bounds, the two separate
error terms, and the scalar test
\begin{equation}\label{calc:scalar-test}
 \frac{L_{m,R,N}(s)}{12(s+3)}>1.
\end{equation}
The uniform inert minorization and its growth-limit theorem make this
a sufficient certificate of $s<\mu$. They also make the test complete:
if $s<\mu$, some $m$ has $A_m^\infty(s)>12(s+3)$, and sufficiently
large finite budgets preserve that strict inequality. This is a
semidecision procedure, with no practical search-length estimate near
$\mu$ and no termination claim when $s\ge\mu$.
The optional \texttt{--search} mode implements the enumeration: at stage
$h$, take $R=N=h$ and test generations $1,\ldots,h$. For any fixed
successful generation the retained sums increase to its full coefficient
along these budgets. The option \texttt{--max-stage} imposes a finite
practical stopping point, whose failure gives no conclusion about $s-\mu$.

For example, with just one allowed real birth and $m=1$, the only
rewarded histories are initial null rings followed by a budding ring.
Writing $a=1/[2(s+2)]$, their exact retained sum is
\[
 L_{1,R,1}(s)=\frac{1-a^R}{2s+3}.
\]
The checker verifies this identity separately from its burst and
individual-clock comparisons.

Run, for example,
\begin{quote}\small\ttfamily
python3 scripts/inert\_growth\_kernel.py --check\\
python3 scripts/inert\_growth\_kernel.py --s 1/2 \textbackslash\\
\hspace*{1em}--generation 2 --core-rings 3 --max-births 4
\end{quote}
The checker compares $180$ small burst calculations against an
independent recursion over individual leaf and core clock events,
including null events. It also checks the potential-profile tail
identities, exact one-ring scalar partial sums, and nested certified
intervals. These are finite verification checks, not a new numerical
exponent bound or a replacement for the infinite coupling proofs.

\section{The uniform newborn growth exponent}
\label{uniform:section}

For $x\in\mathcal X$ put
\[
 F_x(t)=\mathbb E_xC_{\rm tag}(t),\qquad
 G(t)=\sup_xF_x(t),\qquad H(t)=1+\int_0^tG(r)\,dr.
\]
Every $F_x$ is nondecreasing, and $1\le F_x(t)\le e^t$. The same
bounds hold for $G$. At a deterministic or stopping time, partition
the future core genealogy into its currently existing cores and the
families issued from future children of these cores. Each existing
core creates children at intensity at most one, and each newborn
child starts in $\mathcal X$. The compensator formula gives
\begin{equation}\label{uniform:old-core}
 \mathbb E[C(t+u)\mid\mathcal F_t]\le C(t)H(u).
\end{equation}
This holds for any old retained descendant component in an admissible
environment, including the one-stub inert component. It does not
require independence between child families.

In particular, $G(t+u)\le G(t)H(u)$. Also $G(t)\le H(t)$ by the same
decomposition from one newborn core. Exchanging time parameters gives
$G(t+r)\le H(t)G(r)$, whence
\[
 H(t+u)\le H(t)+H(t)\int_0^uG(r)\,dr=H(t)H(u).
\]
Since $G$ is nondecreasing, $H(t)\le1+tG(t)$. These facts prove the
following standard submultiplicative consequence.

\begin{proposition}\label{uniform:lambda}
There is a finite exponent
\begin{equation}\label{uniform:lambda-equation}
 \lambda=\lim_{t\to\infty}\frac{\log G(t)}t
 =\lim_{t\to\infty}\frac{\log H(t)}t
 =\inf_{T>0}\frac{\log H(T)}T,
 \qquad \mu\le\lambda\le1.
\end{equation}
For each $\alpha>\lambda$, a finite constant $K_\alpha$ satisfies
$H(t)\le K_\alpha e^{\alpha t}$ for all $t\ge0$.
\end{proposition}
\begin{proof}
Apply continuous-time subadditivity to the locally bounded function
$\log H$, and use $G\le H\le1+tG$ to identify the two limits. The
inert newborn belongs to $\mathcal X$ and has exponent $\mu$ by
Theorem~\ref{gs:inert-mu}, proving the lower inequality. The upper
inequality follows from $H(t)\le e^t$. For the final assertion choose
$T$ with $H(T)<e^{\alpha T}$, write $t=nT+r$ with $0\le r<T$, and
use submultiplicativity and $H(r)\le e^T$.
\end{proof}

We will first use $\lambda$ without identifying it with $\mu$.
The equality is proved in Section~\ref{erc:closure-section}. This
order is essential: the controlled-moment estimates used in that
proof depend on $\lambda$, and are established independently of the
desired equality.

\section{Uniform comparison of initial environments}

Let $x=(T,u)$ be a finite initial tree with a tagged core vertex $u$.
Only core children born after time zero, and their subsequent descendants,
are assigned positive genealogical generations. Write
\[
 A_m(s;x)=
 \mathbb E_x\sum_{\substack{v\text{ born after }0\\
                          \mathrm{gen}_u(v)=m}}
 e^{-s\tau_v},\qquad m\ge1,\quad s>0.
\]
This is the scalar reward represented by $\mathcal K_s^m\mathbf1$
in Section~\ref{model:section}.
It is not an arbitrary function of the entire ambient tree.

A birth state means that the tag has just been promoted: $L_u=c_u=1$,
so its full degree is two. The original core neighbours of any initial
core vertex remain neighbours forever, and its full degree never
decreases.

\subsection{A common potential genealogy and the comparison identity}

Use a rate-one clock for each actual core vertex. A newly promoted leaf
can be given a fresh independent rate-one core clock, by the Markov
property. Introduce a potential genealogy rooted at $u$ in which
\emph{every} ring of a potential parent's clock creates a potential child,
whether or not that ring creates an actual core child. Label children by
their parent's ring index and run unused potential branches independently.
The actual future core genealogy is a subset of this common potential
genealogy. Its generation counts satisfy
\[
 \mathbb E\overline Z_j(t)=\frac{t^j}{j!}.
\]

Suppose two environments are coupled so that their descendant histories
agree until an increasing event $\Gamma(t)$ occurs. Suppose also that
$\Gamma$ depends only on clock streams of initial core vertices other
than $u$. These streams can be chosen independently of the potential
future genealogy. Each potential generation-$m$ birth has an inclusion
indicator in each model. Their difference has modulus at most one, and
is zero before $\Gamma$. The compensator of potential generation-$m$
births is $\overline Z_{m-1}(t)\,dt$. Hence
\begin{equation}\label{env:comparison}
 \boxed{
 |A_m(s;x)-A_m(s;x')|
 \le
 \int_0^\infty e^{-st}
 \frac{t^{m-1}}{(m-1)!}\,
 \mathbb P(\Gamma(t))\,dt .
 }
\end{equation}
There is no factor two: both actual genealogies are subsets of the
\emph{same} potential genealogy, and the difference of their two
inclusion indicators is bounded by one. Independence is invoked for the
potential event $\Gamma$, not for actual accepted arrows whose rates
depend on the evolving tree.

\subsection{Finite-radius comparison}

Decorate the initial core by the full initial labels $(L_i,c_i)$,
including the numbers of core neighbours outside any displayed ball.

\begin{theorem}[Uniform spatial locality]\label{env:locality}
Suppose two initial environments have isomorphic decorated radius-$R$
balls around their tags, where $R\ge1$. Let $c_u$ be the common initial
core degree of the tag. Then a coupling satisfies
\[
 \mathbb P(\Gamma_R(t))
 \le\min\left\{1,\frac{c_u t^R}{R!}\right\},
\]
and therefore
\begin{equation}\label{env:locality-integral}
 |A_m(s;x)-A_m(s;x')|
 \le
 \int_0^\infty e^{-st}\frac{t^{m-1}}{(m-1)!}
 \min\left\{1,\frac{c_u t^R}{R!}\right\}\,dt .
\end{equation}
For every $T>0$, a convenient explicit bound is
\begin{equation}\label{env:locality-tail}
 \frac{c_u T^{R+m}}{(R+m)R!(m-1)!}
 +s^{-m}e^{-sT}\sum_{j=0}^{m-1}\frac{(sT)^j}{j!}.
\end{equation}
For fixed $m,s$ and a birth state, this tends to zero exponentially in
$R$ on taking $T=aR$ with any fixed $0<a<1/e$.
\end{theorem}
\begin{proof}
For an initial core vertex $i$ with $c_i\ge1$, split its rate-one clock
by first selecting uniformly one of its $c_i$ initial core neighbours.
With probability $c_i/D_i(t)$ attach to that selected neighbour;
otherwise select uniformly among all other current neighbours.
This gives probability $1/D_i(t)$ to every current neighbour and hence
is an exact construction. The potential arrows along each directed
initial core edge $i\to j$ have rate $1/c_i$.

Allow arbitrary discrepancies at the radius-$R$ boundary from time zero.
To affect the tag, any external discrepancy must traverse the unique
initial-core path to the tag by inward arrows in chronological order.
New branches cannot bypass an initial-tree cut edge. Let $\Gamma_R(t)$
be the event that some such path has its required potential arrows
by time $t$.

For a path $u=v_0,v_1,\ldots,v_R$ directed outward from $u$, the expected
number of inward ordered arrow tuples is
\[
 \frac{t^R}{R!}\prod_{j=1}^{R}\frac1{c_{v_j}}.
\]
If $p(\gamma)=\prod_{j=0}^{R-1}c_{v_j}^{-1}$ is the probability that
simple random walk follows that outward path, then
\[
 \prod_{j=1}^{R}\frac1{c_{v_j}}
 =\frac{c_u}{c_{v_R}}p(\gamma)\le c_u p(\gamma).
\]
The probabilities of these length-$R$ outward paths sum to at most one.
A union bound proves the claimed estimate. All source clocks used by
$\Gamma_R$ belong to initial vertices other than $u$, so
\eqref{env:comparison} applies. Cutting the integral at $T$, bounding
$e^{-st}\le1$ on its first part, and evaluating its gamma tail gives
\eqref{env:locality-tail}.
If $c_u=0$, the core is a singleton and there is no exterior to compare.
\end{proof}

The boundary labels are necessary for exact agreement of the interior
initial states. The theorem is an approximation result; it does not say
that the radius-$R$ state itself is Markov.

\subsection{An exact limiting boundary for a high-degree parent}

The absorbing-parent process starts with the tag, its one leaf, and an
inert external core-neighbour stub. The stub contributes one to the tag's
degree. A ring selecting the stub has no effect on the retained process;
it does not kill that process. All remaining dynamics are unchanged.
Denote its scalar generation rewards by $A_m^\infty(s)$.

\begin{theorem}[Uniform parent-degree estimate]\label{env:parent}
If $x=(T,u)$ is a birth state whose parent has initial full degree $d$,
then
\begin{equation}\label{env:parent-bound}
 \boxed{
 |A_m(s;x)-A_m^\infty(s)|
 \le s^{-m}-(s+1/d)^{-m}
 \le\frac{m}{d\,s^{m+1}}.
 }
\end{equation}
\end{theorem}
\begin{proof}
Reserve a potential arrow from the parent to $u$ at constant rate $1/d$.
At such a ring, accept the arrow with probability $d/D_p(t)$.
All rejected and other rings are routed uniformly among the remaining
current neighbours. The accepted rate to $u$ is $1/D_p(t)$ and every
other current neighbour has the same rate, so this construction is exact.
The inequality $D_p(t)\ge d$ makes it valid at every time.

Until the first potential parent arrow, the tag's descendant histories in
the original and absorbing-parent processes agree. The potential arrow
stream is independent of the potential descendant genealogy, and its
probability of ringing by time $t$ is $1-e^{-t/d}$. Substitution in
\eqref{env:comparison} gives the difference of the two gamma integrals
in~\eqref{env:parent-bound}. The final inequality follows by integrating
the derivative of $x^{-m}$ from $s$ to $s+1/d$.
\end{proof}

The limiting process still contains interactions among the tag and its
new descendants. The theorem does not replace that process by independent
families or by a scalar birth-rate model.

\section{A finite approximation of the initial environment}

Fix a birth state, a depth $R\ge1$, and an integer $M\ge3$.
Traverse its initial core outward from the tag. Remove a component at
the first vertex on a path whose initial full degree is at least $M$,
or at depth $R$, whichever is encountered first. Replace each removed
component by an inert core-neighbour stub attached to its retained parent.
Retain all initial leaf counts and all core-neighbour counts of retained
vertices, including their stub neighbours. Write $\pi_{R,M}x$ for this
finite environment with boundary stubs.

Every retained vertex has depth less than $R$ and full degree less than
$M$. Thus, up to rooted isomorphism, the range of $\pi_{R,M}$ is a
finite set of decorated initial environments. Its future evolution is
still an infinite-state process.

\begin{theorem}[Uniform finite initial-environment reduction]
\label{env:finite-initial}
For every birth state $x$, every $m\ge1$, $s>0$, and $T>0$,
\begin{equation}\label{env:finite-bound}
\begin{split}
 |A_m(s;x)-A_m(s;\pi_{R,M}x)|
 \le s^{-m}\left\{
 \frac{T^R}{R!}+\frac{e^T-1}{M}
 +e^{-sT}\sum_{j=0}^{m-1}\frac{(sT)^j}{j!}
 \right\}.
\end{split}
\end{equation}
In particular, for fixed $m,s$, first choosing $T$, then $R,M$,
gives arbitrarily small uniform error over all finite initial environments.
\end{theorem}
\begin{proof}
The depth boundary contributes the potential-event bound $t^R/R!$
from Theorem~\ref{env:locality}, since a birth state has $c_u=1$.

For a first high-degree cut vertex $v$ at distance $r$, reserve its
parentward potential arrow at rate $1/D_v(0)$ and accept with probability
$D_v(0)/D_v(t)$, as in Theorem~\ref{env:parent}.
At intermediate vertices use the old-core potential arrows of rates
$1/c_i$. Cut vertices never occur as intermediate vertices on another
retained cut path, so these constructions are compatible.

If $\gamma=(u=v_0,\ldots,v_r=v)$ is the path to the cut vertex, the
product of inward potential rates is
\[
 \frac1{D_v(0)}\prod_{j=1}^{r-1}\frac1{c_{v_j}}
 =\frac{c_u}{D_v(0)}p(\gamma)
 \le\frac{p(\gamma)}{M}.
\]
First-cut paths are prefix-free, so their simple-random-walk cylinder
probabilities sum to at most one, even when their lengths differ.
The probability of some high-degree influence reaching $u$ by $t$ is
therefore at most
\[
 \frac1M\sup_{r\ge1}\frac{t^r}{r!}
 \le\frac{e^t-1}{M}.
\]
Combining the two kinds of boundary gives a potential event with
probability at most
$\min\{1,t^R/R!+(e^t-1)/M\}$.
Use \eqref{env:comparison}; on $[0,T]$ take its increasing probability
bound at $T$, and use $\int_0^\infty e^{-st}t^{m-1}/(m-1)!\,dt=s^{-m}$.
On $(T,\infty)$ use the probability bound one and the exact gamma tail.
This proves~\eqref{env:finite-bound}.
\end{proof}

The finite range concerns initial environments. Each retained profile still evolves on an infinite state space; the population truncation below supplies a finite calculation at every fixed time.


\subsection{Effective computation at each fixed time}
\label{profile:time}

For rational $T>0$, the preceding pruning construction computes both
$G(T)$ and $H(T)$ with any prescribed rational accuracy. Its exterior
potential event is independent of the future tagged potential
genealogy, and both actual genealogies are subsets of that genealogy.
The potential total core count has mean $e^t$. Consequently, uniformly
for $0\le t\le T$,
\begin{equation}\label{profile:time-error}
 |F_x(t)-F_{\pi_{R,M}x}(t)|
 \le \varepsilon_{R,M}(T)
 :=e^T\left(\frac{T^R}{R!}+\frac{e^T-1}{M}\right).
\end{equation}
If $\mathcal P_{R,M}$ is the finite set of pruned profiles, then
\begin{equation}\label{profile:time-max}
 \max_{P\in\mathcal P_{R,M}}F_P(t)
 \le G(t)\le
 \max_{P\in\mathcal P_{R,M}}F_P(t)+\varepsilon_{R,M}(T).
\end{equation}
The first inequality uses the admissibility of each profile, proved
by finite inert-boundary completion in Section~\ref{model:section}.
Choose $R,M$ effectively to make the displayed error arbitrarily small.

For a fixed profile with $n$ real vertices, enumerate the finite state
space with real population at most $N$, including genealogical marks,
and kill on the next birth. All transition rates are rational. If
$Y_n(t)$ is a Yule process started at $n$, the omitted mean is at most
\begin{equation}\label{profile:population-error}
 \mathbb E[Y_n(t);Y_n(t)>N]
 \le\frac{\mathbb EY_n(t)^2}{N}
 =\frac{n(n+1)e^{2t}-ne^t}{N}.
\end{equation}
Indeed, the real birth rate is at most the population; inert choices
are null. Couple with $Y_n$ so that a crossing of the retained cutoff
forces $Y_n(t)>N$, and use $C_{\rm tag}(t)\le Y_n(t)$.

A rational uniformization rate exceeding every finite-state exit rate
expresses its matrix exponential as a Poisson-weighted series of
powers of a rational substochastic matrix. Its partial sums and a
certified Poisson tail give rational enclosures. Elementary rational
enclosures for exponentials make \eqref{profile:population-error}
effective as well.

Finally every $F_x$ is $e^T$-Lipschitz on $[0,T]$: its core-birth
intensity is at most $F_x(t)\le e^t$. The supremum $G$ has the same
Lipschitz constant. Rational-grid quadrature, combined with
\eqref{profile:time-max} and the population cutoff, computes
$H(T)=1+\int_0^TG(t)\,dt$. This argument concerns a fixed time and
uses no convergence rate for $t^{-1}\log G(t)$.

\section{Finite-window supersolutions with bounded increments}
\label{sc:section}

Let $P_t$ denote the inert-boundary semigroup.  An ordinary vertex
addition means a single leaf birth, which changes the core count by
either zero or one.
Write $G(t)$ for the supremum of newborn retained means over finite
environments and $H(t)=1+\int_0^tG(r)\,dr$.  By Section~\ref{uniform:section},
\[
 H(t+u)\le H(t)H(u),\qquad
 \lambda=\lim_{t\to\infty}\frac{\log H(t)}t,
 \qquad \mathbb E_y^\infty C(t)\le C(y)H(t).
\]
Each $H(T)$ for rational $T>0$ is effectively computable by the certified
finite-profile and finite-population reductions.

\begin{lemma}[Finite-time sensitivity]
\label{sc:sensitivity}
Let $\psi$ be any function with ordinary-birth increments bounded in
absolute value by $J$.  Assuming the expectations below exist, for every
ordinary vertex addition $y\mapsto y'$ and every $t\ge0$,
\begin{equation}\label{sc:sensitivity-bound}
 |P_t\psi(y')-P_t\psi(y)|\le J e^{2t}.
\end{equation}
\end{lemma}
\begin{proof}
Couple the two full vertex processes using common rate-one clocks at
common vertices and maximal couplings of their neighbour choices.
Initially all vertices of $y$ are common, and the added vertex is
unmatched.  Match newly born vertices only when their source event and
target agree in both copies; otherwise use distinct new labels.
The common vertices consequently induce a connected tree containing the
whole initial $y$.  Each component of unmatched vertices attaches to
this common tree through exactly one edge.

Let $Z(t)$ be the total number of unmatched vertices across both copies.
An unmatched vertex clock rings at rate one and creates at most one
new unmatched vertex.  At a common vertex, the probability of a
neighbour-choice mismatch is
\[
 \frac{\max(a,b)}{d_0+\max(a,b)}\le\frac{a+b}{2},
\]
where $d_0\ge1$ is the number of shared neighbours (counting the shared
inert stub at the root) and $a,b$ are the two counts of unmatched
neighbours.  The inequality follows directly for the integer cases
$a+b=1$ and $a+b\ge2$, and is trivial when $a+b=0$.
Summing over common vertices bounds the mismatch rate by half the
number of common--unmatched edges, and hence by $Z(t)/2$.
Each such mismatch creates at most two unmatched vertices.  Inert-stub
selections can only reduce these increments.  Hence
\[
 \frac{d}{dt}\mathbb EZ(t)\le2\mathbb EZ(t),
 \qquad Z(0)=1,
\]
so $\mathbb EZ(t)\le e^{2t}$.

At any time the two trees can each be built from their common connected
tree by ordinary leaf additions, with total number of additions $Z(t)$.
The common tree still contains the initial rooted core and inert stub,
so all those intermediate states are legal.  Telescoping the jump bound
gives $|\psi(X_t')-\psi(X_t)|\le JZ(t)$.  Taking expectations proves
\eqref{sc:sensitivity-bound}.
\end{proof}

\begin{theorem}[A finite time-window supersolution]
\label{sc:window}
Let $s,T>0$, and suppose a certified rational $q<1$ satisfies
\begin{equation}\label{sc:time-test}
 e^{-sT}H(T)\le q.
\end{equation}
Define
\begin{equation}\label{sc:window-U}
 \boxed{U_{s,T}(y)=\frac1{1-q}
                \int_0^T e^{-st}\mathbb E_y^\infty C(t)\,dt.}
\end{equation}
Then $U_{s,T}$ is nonnegative and, for computable $s,T$, effectively
computable at every finite state.  It satisfies
\begin{equation}\label{sc:window-properties}
 (s-\mathcal L_\infty)U_{s,T}\ge C,
 \qquad
 |U_{s,T}(y')-U_{s,T}(y)|
 \le B_{s,T}:=\frac1{1-q}\int_0^T e^{(2-s)t}\,dt
\end{equation}
for every ordinary vertex addition.  The final integral equals
$(e^{(2-s)T}-1)/(2-s)$ when $s\ne2$, and $T$ when $s=2$.
Moreover
\begin{equation}\label{sc:window-linear}
 0\le U_{s,T}(y)\le a_{s,T}C(y),\qquad
 a_{s,T}:=\frac1{1-q}\int_0^T e^{-st}H(t)\,dt.
\end{equation}
The elementary bound $H(t)\le e^t$ permits the explicit replacement
$a_{s,T}\le\int_0^T e^{(1-s)t}dt/(1-q)$ whenever only an upper
constant is needed.
\end{theorem}
\begin{proof}
The finite-horizon resolvent identity gives
\[
 (s-\mathcal L_\infty)
      \int_0^T e^{-st}P_tC\,dt
 =C-e^{-sT}P_TC.
\]
It is justified pointwise by finite state localization and Yule
domination.  The conditional old-core mean bound gives
$P_TC(y)\le C(y)H(T)$.
The right side of the resolvent identity is consequently at least
$(1-q)C$, proving the supersolution inequality.
Lemma~\ref{sc:sensitivity} applied to $\psi=C$, $J=1$ proves the
increment estimate after integration.  The old-core mean bound also
gives \eqref{sc:window-linear}.

At a fixed finite initial state, rational finite-state generator
truncations compute $P_tC$ uniformly over $[0,T]$: the real population
is dominated by a Yule process, and $C\le N$ gives an explicit
discarded-population error.  Uniformization supplies rational intervals
for the matrix-exponential entries.  Time quadrature is effective as
well, since $0\le\mathcal L_\infty C\le N$ and
$\mathbb E N(t)\le N(0)e^t$ bound its finite-horizon derivative.
Thus the integral defining $U_{s,T}$ is computable with certified error.
\end{proof}

\begin{corollary}[Completeness above $\lambda$]
\label{sc:complete-lambda}
For every rational $s>\lambda$, a search over positive integer $T$ and
certified upper approximations of $H(T)$ eventually verifies
\eqref{sc:time-test} with $q=1/2$, and hence constructs a supersolution
with a uniform
ordinary-birth increment at $s$.
\end{corollary}
\begin{proof}
Since $s>\lambda$, $e^{-sT}H(T)\to0$.  Thus the strict inequality in
\eqref{sc:time-test} holds for sufficiently large integer $T$ with
$q=1/2$.  Dovetailing the positive integers $T$ and approximation
precisions eventually certifies that strict inequality.  The preceding
theorem then supplies the explicit function and increment bound.
\end{proof}


\section{Predictable boundary inputs preserve the uniform moment exponents}
\label{cm:section}

Let $\mathcal Q$ be the class of nonanticipating input policies which add
a leaf at the retained root at predictable intensity $q_t\in[0,1/2]$.
Apart from these additions, the component follows the inert generator.
The root starts in the newborn state $x_\infty$, with two real vertices.
The policies may use the observed component history and additional
randomness. An actual exterior supplies a policy in this class.

Theorem~\ref{sc:window} implies that,
for every $\alpha>\lambda$, there is a function $U$ and finite positive
constants $m,A,J$ such that
\begin{equation}\label{cm:window-input}
 mC\le U\le AC,\qquad
 (\alpha-\mathcal L_\infty)U\ge C,\qquad
 |U(y')-U(y)|\le J
\end{equation}
for every ordinary leaf addition $y\to y'$. In particular, one may use
$U_{\alpha,T}$ from that construction, with
$m=(1-e^{-\alpha T})/[\alpha(1-q)]$ and its already stated $A,J$.
Here the scalar $q<1$ is the finite-window contraction parameter.

\begin{theorem}[All controlled moments have the same normalized rate]
\label{cm:all-moments}
For every real $p\ge1$ and every $a>\lambda$, there is a finite constant
$K_{p,a}$, independent of the input policy, such that
\begin{equation}\label{cm:bound}
 \sup_{q\in\mathcal Q}\mathbb E_q C(t)^p
 \le K_{p,a}e^{pa t}\qquad(t\ge0).
\end{equation}
Consequently,
\begin{equation}\label{cm:exponent}
 \boxed{\lim_{t\to\infty}\frac1t
 \log\sup_{q\in\mathcal Q}\mathbb E_q C(t)^p=p\lambda.}
\end{equation}
For computable parameters and certified positive slack, the constants
are effective once a finite-window certificate at any
$\alpha\in(\lambda,a)$ has been supplied. More generally, the same
moment upper bound holds from any fixed finite initial component, with
constants depending on that component.
\end{theorem}
\begin{proof}
Choose $\lambda<\alpha<a$. The established lower bound on $\mu$ gives
$\alpha>\lambda\ge\mu\ge0.565>1/2$. Set
\[
 W=U+\frac{J}{2\alpha}.
\]
For every admissible policy, its conditional generator satisfies
\begin{equation}\label{cm:linear-drift}
 (\mathcal L_\infty+q_t\Delta_{\rm root})W
 \le\alpha U-C+J/2\le\alpha W.
\end{equation}
The jumps of $W$ have absolute size at most $J$, and $W\ge mC\ge m$.

Let $N_t$ be the real vertex count and put $Y_t=N_t+1/2$. Its birth
intensity is at most $Y_t$. For each integer $r\ge1$, the rising
factorial $y^{\overline r}=y(y+1)\cdots(y+r-1)$ satisfies
\[
 y\bigl[(y+1)^{\overline r}-y^{\overline r}\bigr]
 =r y^{\overline r}.
\]
It follows by localization that
\begin{equation}\label{cm:population}
 \mathbb E Y_t^r\le
 \mathbb E Y_t^{\overline r}
 \le Y_0^{\overline r}e^{rt}.
\end{equation}

First take an even integer $p=2r\ge2$. Taylor's formula and
$|\Delta W|\le J$ give, at every transition,
\[
 (W+\Delta W)^p-W^p-pW^{p-1}\Delta W
 \le\frac{p(p-1)}2J^2(W+J)^{p-2}.
\]
Since $W\ge m$, set
$K_p=\tfrac12p(p-1)J^2(1+J/m)^{p-2}$. Summing over transitions and
using \eqref{cm:linear-drift} yields
\[
 (\mathcal L_\infty+q_t\Delta_{\rm root})W^p
 \le p\alpha W^p+K_p W^{p-2}Y.
\]
For $p>2$, Young's inequality, or maximization in $W$, gives for every
$\eta>0$ a finite explicit $B_{p,\eta}$ with
\[
 K_p W^{p-2}Y\le\eta W^p+B_{p,\eta}Y^{p/2}.
\]
For example, one may take
\[
 B_{p,\eta}=\frac2p K_p
 \left(\frac{(p-2)K_p}{p\eta}\right)^{(p-2)/2}.
\]
For $p=2$, the corresponding remainder is simply $J^2Y$.
Choose $\eta=p(a-\alpha)$ when $p>2$; in the $p=2$ case increase
the drift coefficient $2\alpha$ harmlessly to $2a$. Thus
\[
 \frac d{dt}\mathbb E W_t^p
 \le pa\mathbb E W_t^p+B\,Y_0^{\overline {p/2}}e^{pt/2},
\]
with $B=B_{p,\eta}$ for $p>2$ and $B=J^2$ for $p=2$.
Because $a>1/2$, integration gives
\begin{equation}\label{cm:explicit-even}
 \mathbb E W_t^p
 \le e^{pa t}\left[
 W_0^p+\frac{B\,Y_0^{\overline {p/2}}}{p(a-1/2)}\right].
\end{equation}
Divide by $m^p$. All localization limits are justified by
$W\le AC+J/(2\alpha)\le AN+J/(2\alpha)$ and the finite-horizon
population moments. These estimates hold uniformly for every policy.

For a general real $p\ge1$, choose an even integer $k\ge p$ and use
$\mathbb E C^p\le(\mathbb E C^k)^{p/k}$. This proves
\eqref{cm:bound}. Conversely the controlled class contains every admissible
newborn process, including inert exterior profiles. Jensen's inequality gives
\[
 \sup_{q\in\mathcal Q}\mathbb E_q C(t)^p
 \ge \left(\sup_x\mathbb E_x C_{\rm tag}(t)\right)^p=G(t)^p.
\]
The established limit $t^{-1}\log G(t)\to\lambda$, followed by
$a\downarrow\lambda$ in the upper bound, proves
\eqref{cm:exponent}.
\end{proof}


\section{A uniform freeze-or-grow bound for the external parent}
\label{ps:section}

The statement below applies to the original core process, also allowing
inert boundary stubs.  A stub is counted in its incident core's degree,
receives outward selections without changing state, and sends no arrivals.
No monotonicity of expected descendant counts is used.

\begin{lemma}[Reciprocal population under persistent immigration]
\label{ps:reciprocal}
Let $0<\rho\le1$, and let $X$ be the birth--death chain on
$\mathbb Z_{\ge0}$ with birth rate $x+\rho$ at $x$ and death rate one
at every $x\ge1$.  If $X(0)=0$, then
\begin{equation}\label{ps:reciprocal-bound}
 \mathbb E\frac1{X(t)+1}
 \le \left(1+\frac3\rho\right)e^{-\rho t/3}
 \qquad(t\ge0).
\end{equation}
\end{lemma}
\begin{proof}
Put $q_x=\Pr\{\operatorname{Pois}(1)\ge x\}$ and
\[
 f(x)=\frac1{x+1},\qquad
 h(x)=f(x)+\frac3\rho q_x.
\]
The generator $G_0$ with birth rate $x$ and unit death rate satisfies
$G_0q_x=0$ for $x\ge1$, because
\[
 q_{x+1}-q_x=-\frac{e^{-1}}{x!}.
\]
At $x=0$, $G_0q_0=0$ as well.  For $x\ge1$,
\[
 G_0f(x)=-\frac{x-2}{x(x+2)}.
\]
Consequently, for the generator $G$ including immigration,
\[
 Gh(x)=G_0f(x)
       -\frac{\rho}{(x+1)(x+2)}-\frac3{e\,x!}
 \quad(x\ge1),
\]
whereas $Gh(0)=-\rho/2-3/e$.

For $x\ge1$, summing the factorial tail gives
$q_x\le2/(e\,x!)$.  For $x\ge4$ we also have
$G_0f(x)\le-(5/12)f(x)$.  Since $\rho\le1$, these inequalities imply
\[
 Gh(x)+\frac\rho3 h(x)
 \le-\frac1{12}f(x)-\frac1{e\,x!}<0
 \qquad(x\ge4).
\]
At $x=0,1,2,3$, respectively, the exact expressions for
$Gh(x)+(\rho/3)h(x)$ are
\[
 -\frac\rho6+1-\frac3e,\quad
 \frac43-\frac4e,\quad
 \frac\rho{36}+1-\frac7{2e},\quad
 -\frac1{15}+\frac\rho{30}+1-\frac3e.
\]
They are all negative, using $e<3$ and $\rho\le1$.
Thus $Gh\le-(\rho/3)h$ at every state.

The chain is nonexplosive, and $h$ is bounded.  Localization and
Dynkin's formula therefore yield
$\mathbb E h(X(t))\le h(0)e^{-\rho t/3}$.
Since $f\le h$ and $h(0)=1+3/\rho$, this proves the claim.
\end{proof}

Fix a real core vertex $p$ whose initial full degree is at most an
integer $M\ge2$, and set
\[
 \tau_M=\inf\{t:D_p(t)>M\},\qquad
 C_M=1+3M,\qquad \kappa_M=\frac1{3M}.
\]
For real neighbours $v$ of $p$, define
\[
 \lambda_p(t)=\sum_{v\sim p,\ v\text{ real}}\frac1{D_v(t)}.
\]
Thus $\lambda_p$ is the rate at which $p$ receives a new leaf and
its full degree increases.  Its leaf-promotion rate is
$b_p=L_p/D_p$.

\begin{theorem}[Uniform local stabilization]
\label{ps:stabilization}
Uniformly over all finite initial environments and their boundary stubs,
\begin{equation}\label{ps:rate-bound}
 \mathbb E\bigl[\mathbf1_{\{t<\tau_M\}}\lambda_p(t)\bigr]
 \le M C_M e^{-\kappa_Mt}
       \sum_{j=0}^{M-1}\frac{(C_Mt)^j}{j!}.
\end{equation}
In particular, the probability that $p$ has a degree increase or a leaf
promotion at some time in $(T,\tau_M]$ is at most
\begin{align}
 Q_M(T)
 &=\min\left\{1,
 2M C_M\sum_{j=0}^{M-1}\frac{C_M^j}{j!}
       \int_T^\infty e^{-\kappa_Mt}t^j\,dt\right\}
 \label{ps:tail-integral}\\
 &=\min\left\{1,
 2M C_M e^{-\kappa_MT}
 \sum_{j=0}^{M-1}C_M^j
       \sum_{i=0}^j\frac{T^i}{i!\,\kappa_M^{j-i+1}}
 \right\}.
 \label{ps:tail-explicit}
\end{align}
This bound tends to zero exponentially in $T$ for each fixed $M$.
\end{theorem}
\begin{proof}
While $D_p\le M$, implement the rate-one clock at $p$ using $M$
independent slots, each of rate $1/M$.  Give every neighbour a permanent
distinct slot when that neighbour first appears.  At an occupied slot,
choose its assigned neighbour.  At an empty slot, choose uniformly
among the current $D_p$ neighbours.  Every particular neighbour is
then chosen at rate
\[
 \frac1M+\frac{M-D_p}{M D_p}=\frac1{D_p},
\]
so the construction is exact until $\tau_M$.

Let $v$ be a real neighbour present initially, or born as a new
neighbour of $p$ at a stopping time $\sigma<\tau_M$.  Its occupied
slot supplies arrivals at rate $1/M$.  Every such arrival creates a
leaf at $v$; the first one also turns $v$ into a core if it was still
a leaf.  The number of leaf neighbours of $v$ lower-dominates, until
$\tau_M$, the chain of Lemma~\ref{ps:reciprocal} with $\rho=1/M$
and initial state zero.  Use the same leaf-birth Poisson measure,
the occupied slot for immigration, and a death at every ring of
$v$'s own rate-one clock.  Actual leaf losses occur at most once at
such a ring, and all additional arrivals only increase the actual
leaf count.  These three driving processes are independent before
stopping.  They can be continued after stopping to define the full
comparison chain.

The permanent edge $pv$ gives $D_v\ge X+1$ before stopping.  Therefore,
conditionally on the history at $\sigma$, for every $a\ge0$,
\begin{equation}\label{ps:conditional-kernel}
 \mathbb E\left[
   \frac{\mathbf1_{\{\sigma+a<\tau_M\}}}{D_v(\sigma+a)}
   \,\middle|\,\mathcal F_\sigma\right]
 \le C_Me^{-\kappa_Ma}.
\end{equation}
Initial neighbours use $\sigma=0$.  Stubs contribute zero and require
no comparison chain.

Give each new neighbour $w$ of $p$ a genealogical parent: the old
neighbour $v$ whose clock selected $p$ and created $w$.
This is a genealogy of neighbours, separate from the core genealogy.
Its roots are the at most $M$ initial real neighbours.  Each individual's
conditional mean birth intensity, as a function of its age before
stopping, is bounded by the kernel
$k_M(a)=C_Me^{-\kappa_Ma}$ in
\eqref{ps:conditional-kernel}.
Iterated conditional first moments, or the compensator formula followed
by Tonelli's theorem, bound the expected birth intensity contributed
by neighbour generation $j$ by
\[
 M k_M^{*(j+1)}(t)
 =M C_M^{j+1}e^{-\kappa_Mt}\frac{t^j}{j!}.
\]
No independence between neighbours is asserted or needed.
Before $\tau_M$ there are at most $M$ neighbours, so no contributing
individual can have neighbour generation greater than $M-1$.
Summing gives~\eqref{ps:rate-bound}.

Finally, each leaf neighbour contributes one to $\lambda_p$, so
$b_p=L_p/D_p\le L_p\le\lambda_p$.  The total intensity of changes
to $(D_p,L_p,c_p)$ is thus at most $2\lambda_p$.
The compensator bound for the number of such changes after $T$, counting
also the degree increase at $\tau_M$, proves
\eqref{ps:tail-integral}.  Integrating the integer powers gives
\eqref{ps:tail-explicit}.
\end{proof}


\section{Polynomial likelihood reduction of a bounded-degree parent}
\label{plc:section}

Let $Y_t$ be the retained descendant component of a newborn tag, viewed
with its parent edge replaced by an inert stub.  Let $D_t$ be the full
degree of its real external parent.  Fix $M\ge2$ and an initial degree
$d\in\{2,\ldots,M\}$, and stop both coordinates at
\[
 \tau_M=\inf\{t:D_t=M+1\}.
\]
Before this stopping time, the parent's degree-increase intensity is
\[
 \lambda_p(t)=\sum_{v\sim p}\frac1{D_v(t)}\le D_t\le M.
\]
The sum includes real leaf neighbours, which contribute one, and
inert neighbours, which contribute zero.  Since the tag is a real
neighbour, $\lambda_p(t)>0$ at every finite state.
There are at most $K=M+1-d$ parent degree increases before stopping.

\begin{theorem}[A finite family of reference laws with polynomial loss]
\label{plc:theorem}
There is a reference law $\mathbb Q_d$ for $(Y,D)$, depending only on
$d$ and $M$, with time-dependent generator before stopping
\begin{align}
 \mathcal G_t^{\rm ref}f(y,d)
 &=\mathcal L_\infty f(\cdot,d)(y)
   +\frac1d\{f(y^+,d)-f(y,d)\}\notag\\
 &\quad+\frac1{(1+t)^2}\{f(y,d+1)-f(y,d)\},
 \label{plc:reference-generator}
\end{align}
such that every actual finite exterior with parent degree $d$ satisfies
the following bound.  For every $T\ge0$ and every nonnegative measurable
functional $Z$ of the stopped retained-and-degree path through time $T$,
\begin{equation}\label{plc:domination}
 \mathbb E_{
m actual} Z
 \le e\,[M(1+T)^2]^K\,\mathbb E_{\mathbb Q_d} Z.
\end{equation}
Thus arbitrary exterior histories with a bounded-degree parent can be
replaced by one of finitely many explicit reference models at a
polynomial, rather than exponential, cost in the time horizon.
\end{theorem}
\begin{proof}
Use the full finite-tree birth process until $\tau_M$.  A parent degree
increase is exactly one leaf addition at the parent, with total rate
$\lambda_p(t)$.  Change this single transition channel to the rate
$h(t)=(1+t)^{-2}$ and leave every other transition rate unchanged.
Both the original and the modified processes are nonexplosive: the
modified total real birth rate is at most the original population plus
one.  The usual finite-horizon counting-process change of measure gives
the actual law relative to the modified law with density
\[
 Z_T^{\rm RN}
 =\exp\!\left\{\int_0^{T\wedge\tau_M}
                     (h(t)-\lambda_p(t))\,dt\right\}
   \prod_{j:\,\sigma_j\le T\wedge\tau_M}
        \frac{\lambda_p(\sigma_j-)}{h(\sigma_j)}.
\]
Here $\sigma_j$ are the parent degree-increase times, including the
stopping increase when it occurs before $T$.  Localization in total
population first reduces the calculation to bounded finite-state
intensities.  Alternatively, the density follows directly by taking
the ratio of the finite jump-history densities and then passing to the
nonexplosive limit.

There are at most $K$ factors in the product,
$\lambda_p\le M$, and $\int_0^\infty h(t)dt=1$.  Hence pathwise
\[
 Z_T^{\rm RN}\le e\,[M(1+T)^2]^K.
\]

It remains to identify the marginal law of the retained coordinates
under the modified full process.  A parent degree increase changes only
$D$, and now has deterministic hazard $h(t)$.  A parent selection of the
tag creates one retained root leaf at rate $1/D$.  All other retained
transitions have exactly the inert generator $\mathcal L_\infty$.
Every remaining exterior transition changes neither $Y$ nor $D$.
Thus the pair $(Y,D)$ has precisely the autonomous generator
\eqref{plc:reference-generator}; in particular its law no longer depends
on the rest of the initial exterior.  This proves
\eqref{plc:domination}.
\end{proof}

The reference degree chain is driven by an independent inhomogeneous
Poisson clock with integrated hazard one, stopped after $K$ jumps.
The input $1/D_t$ is nonincreasing, changes at most $K$ times, and is
eventually constant on every path that does not reach the stopping
state.  For example, taking
$Z=C(Y_T)\mathbf1_{\{T<\tau_M\}}$ gives a uniform bound for the part
of the tagged mean that remains behind a low-degree parent.

More generally one may use
$h(t)=\eta(1+t)^{-\beta}$ with $\eta>0$ and $\beta>1$.
The likelihood factor is then at most
\[
 \exp\!\left(\frac{\eta}{\beta-1}\right)
 \max\left\{1,\frac M\eta(1+T)^\beta\right\}^K.
\]
Thus the total reference jump intensity can be made arbitrarily small
without introducing an exponential loss in~$T$.

\subsection{An exponential occupation bound after the root is large}

\begin{proposition}[Uniform inverse-degree exponential moment]
\label{plc:occupation}
Suppose a retained root starts with full degree $D_0\ge H\ge2$ and
receives incoming leaves at predictable intensity at least
$\rho\in(0,1]$.  The other incoming events may be arbitrary events of
the usual core process.  Put
\[
 c_\rho=1+\frac3\rho,\qquad
 \kappa_\rho=\frac\rho3,\qquad
 A_{H,\rho}=\frac{\log(c_\rho H)+1}{\kappa_\rho H}.
\]
Then
\begin{equation}\label{plc:occupation-mean}
 \mathbb E\int_0^\infty\frac{dt}{D_t}\le A_{H,\rho}.
\end{equation}
For every $p>0$ with $pA_{H,\rho}<1$,
\begin{equation}\label{plc:occupation-exponential}
 \boxed{
 \mathbb E\exp\!\left(p\int_0^\infty\frac{dt}{D_t}\right)
 \le\frac1{1-pA_{H,\rho}}.
 }
\end{equation}
The estimates are uniform over all finite internal environments, initial
leaf counts and other permitted inputs.  They also hold with the
integrals stopped at any stopping time up to which the lower immigration
bound remains valid.
\end{proposition}
\begin{proof}
The root degree is nondecreasing, so $D_t\ge H$.  Lower-couple its
leaf count with the birth--death chain $X$ started at zero, with births
at rate $x+\rho$ and deaths at rate one whenever $x>0$.  A rate-$\rho$
stream is extracted from the incoming leaves, the leaf births are
coupled by their Poisson random measures, and the comparison loses
one leaf at every root clock ring.  Actual root leaf losses occur at
most once at such a ring.  This is the same comparison used in
Lemma~\ref{ps:reciprocal}.  Its reciprocal estimate gives
\[
 \mathbb E\frac1{D_t}
 \le\mathbb E\min\left\{\frac1H,\frac1{X_t+1}\right\}
 \le\min\left\{\frac1H,c_\rho e^{-\kappa_\rho t}\right\}.
\]
Integrating the last bound, with the two terms equal at
$t=\log(c_\rho H)/\kappa_\rho$, proves
\eqref{plc:occupation-mean}.

The same bound applies conditionally after every stopping time: the
root degree is still at least $H$, the comparison chain may be restarted
at zero below its actual leaf count, and the remaining input still has
intensity at least $\rho$.  For
$I=\int_0^\infty D_t^{-1}dt$, iterated conditioning in the ordered
time integrals therefore gives the Khasminskii bounds
\[
 \mathbb E I^n\le n!A_{H,\rho}^n\qquad(n\ge0).
\]
Summing the nonnegative exponential series proves
\eqref{plc:occupation-exponential}.  With a stopping time, insert the
indicator that the input lower bound is still in force; the same
conditional argument only decreases the integrals.
\end{proof}

In particular, $A_{H,\rho}\to0$ as $H\to\infty$ for fixed $\rho$.
Thus any prescribed exponential moment is uniformly finite once the
root degree is sufficiently large.


\section{The bounded-parent survival gap}
\label{pup:section}

By Theorem~\ref{cm:all-moments}, for every $\alpha>\lambda$ and every fixed $r>1$,
\begin{equation}\label{pup:controlled-second}
 \sup_{q\in\mathcal Q}\mathbb E_q C(t)^r\le K_{r,\alpha}e^{r\alpha t}.
\end{equation}
The constants are uniform in the newborn exterior. This estimate does not use $\lambda=\mu$.

\subsection{Uniform control while the actual parent remains bounded}

Fix $M\ge2$ and, for a newborn tag whose parent starts at degree
$d\le M$, let $\tau_M=\inf\{t:D_p(t)=M+1\}$.  Define
\[
 L_M(t)=\sup_{x:\,D_p(0)\le M}
       \mathbb E_x[C_{\rm tag}(t)\mathbf1_{\{t<\tau_M\}}].
\]
The supremum is over the actual finite exteriors, with the same permitted
inert stubs as in the earlier profile construction.

\begin{lemma}[A stopped inverse-likelihood moment]\label{pup:stopped-I}
Under the reference law replacing parent degree gains by
$h(t)=(1+t)^{-2}$, stopped at $\tau_M$, put
\[
 I_M=\int_0^{\tau_M}\lambda_p(t)\,dt,
 \qquad B_M=3M^2(1+3M).
\]
Then, for $aB_M<1$,
\begin{equation}\label{pup:stopped-moment}
 \mathbb E e^{aI_M}\le(1-aB_M)^{-1}.
\end{equation}
The expected remaining occupation is at most $B_M$ conditionally after
every stopping time before $\tau_M$.
\end{lemma}
\begin{proof}
Before stopping there are at most $M$ neighbour labels of $p$ in total.
Each real neighbour, from its time of birth, receives immigration at
rate at least $1/M$, by the permanent-slot construction.  The reciprocal
population bound gives a conditional expected remaining contribution
of at most
\[
 \int_0^\infty(1+3M)e^{-t/(3M)}dt=3M(1+3M)
\]
to the integral of its reciprocal degree.  The estimate is valid
also for a neighbour initially of degree one and for stopped integrals.
Summing over at most $M$ present or future neighbour labels gives $B_M$.
The same estimate restarts conditionally from any stopping time; the
number of remaining possible neighbour labels is still at most $M$.
Iterated ordered-time integration gives
$\mathbb E I_M^r\le r!B_M^r$, and summation proves
\eqref{pup:stopped-moment}.
\end{proof}

\begin{theorem}[The bounded-parent survival gap]\label{pup:low-survival}
If $\mu<\lambda$, then for every fixed integer $M\ge2$ there is
$\gamma_M<\lambda$ such that
\begin{equation}\label{pup:low-survival-bound}
 \limsup_{t\to\infty}\frac1t\log L_M(t)\le\gamma_M.
\end{equation}
Thus an exponent gap cannot be supported by newborn histories whose
parent remains below a fixed degree threshold throughout observation.
\end{theorem}
\begin{proof}
Fix $\varepsilon\in(0,1)$ and split the event $t<\tau_M$ according
to whether a parent degree gain occurs after time $\varepsilon t$.
Theorem~\ref{ps:stabilization} gives
\begin{gather*}
 \Pr\{\text{such a late gain before or at }\tau_M\}
 \le Q_M(\varepsilon t),\\
 Q_M(r)=O_M(\operatorname{poly}(r)e^{-\kappa_M r}),
 \qquad \kappa_M=1/(3M).
\end{gather*}
By \eqref{pup:controlled-second}, the contribution of these histories
has upper exponent at most $\alpha-\kappa_M\varepsilon/2$ for every
$\alpha>\lambda$.

Let $A_t$ be the other event: $t<\tau_M$ and all parent degree gains
through time $t$ occur by $\varepsilon t$.  Theorem~\ref{plc:theorem} bounds the actual expectation on $A_t$ by a fixed polynomial
in $t$ times the reference expectation, uniformly in the initial
exterior.  Only the finitely many starting degrees $d=2,\ldots,M$
need be considered.

Realize each reference marginal by modifying a fixed genuine finite
star completion of the newborn tag.  Let $Z_t$ be its actual-to-reference
likelihood, and let $N(r)$ be its full reference real population, stopped
at $\tau_M$.  Before stopping,
\[
 \lambda_p(r)\ge1/D_u(r)\ge1/N(r),\qquad h(r)\le1.
\]
There are at most $M$ parent degree gains.  Hence on $A_t$,
\[
 Z_t^{-1}
 =\exp\!\left(\int_0^t(\lambda_p-h)dr\right)
       \prod_{\sigma_j\le t}\frac{h(\sigma_j)}{\lambda_p(\sigma_j-)}
 \le e^{I_M}N(\varepsilon t)^M.
\]
Choose $a>0$ with $2aB_M<1$.  By Cauchy--Schwarz and
Lemma~\ref{pup:stopped-I},
\begin{equation}\label{pup:prefix-inverse}
 \mathbb E_{\rm ref}[Z_t^{-a}\mathbf1_{A_t}]
 \le K_{a,M}e^{aM\varepsilon t}.
\end{equation}
Indeed, the reference total birth rate is at most $N+1$; a shifted Yule
domination gives $\mathbb E N(r)^v\le K_{M,v}e^{vr}$ for every fixed
$v>0$.

Put $\theta=a/(a+2)$, so $a=2\theta/(1-\theta)$.  Three-factor
H\"older now gives
\begin{align*}
 \mathbb E_{\rm ref}[C(t)\mathbf1_{A_t}]
 &\le
 \bigl(\mathbb E_{\rm ref}[Z_tC(t)\mathbf1_{A_t}]\bigr)^\theta
 \bigl(\mathbb E_{\rm ref}C(t)^2\bigr)^{(1-\theta)/2}\\
 &\qquad\times
 \bigl(\mathbb E_{\rm ref}[Z_t^{-a}\mathbf1_{A_t}]\bigr)^{(1-\theta)/2}.
\end{align*}
The first expectation is at most the total core mean from the fixed
genuine star completion, and therefore has upper exponent $\mu$.
The second has upper exponent $2\alpha$ by the controlled-moment
bound; stopping only decreases its core count.  Formula
\eqref{pup:prefix-inverse} bounds the third.  After the polynomial
forward likelihood factor, the good contribution consequently has
upper exponent
\[
 \theta\mu+(1-\theta)\alpha+\theta M\varepsilon.
\]
Take $a=1/(4B_M)$, so $\theta=1/(8B_M+1)$, and write
$g=\lambda-\mu>0$.  Choose $\varepsilon=g/(4M)$ and then
$\alpha=\lambda+\delta$ with
\[
 0<\delta<
 \min\left\{\frac{\theta g}{4(1-\theta)},
              \frac{\kappa_M\varepsilon}{4}\right\}.
\]
Both displayed contribution exponents are strictly less than $\lambda$:
the good one is below $\lambda-\theta g/2$, and the bad one below
$\lambda-\kappa_M\varepsilon/4$.  The finite maximum over $d$ changes
only constants.  This proves \eqref{pup:low-survival-bound}.
\end{proof}

\begin{corollary}[A bounded retained root is also subcritical]
\label{pup:root-stay}
Under $\mu<\lambda$, for every fixed integer $H\ge2$,
\[
 \limsup_{t\to\infty}\frac1t\log\sup_x
 \mathbb E_x[C_{\rm tag}(t)\mathbf1_{\{D_u(t)\le H\}}]<\lambda.
\]
\end{corollary}
\begin{proof}
Partition the tagged core genealogy into the root and the descendant
components of its future children. On $D_u(t)\le H$, each child's
parent has stayed at degree at most $H$ since its birth, because degree
is nondecreasing. Conditional on the child's birth, its weighted
remaining mean is therefore at most $L_H$ at the remaining horizon.
The root produces children at intensity at most one. The compensator
formula and positivity give
\[
 \sup_x\mathbb E_x[C_{\rm tag}(t)\mathbf1_{\{D_u(t)\le H\}}]
 \le1+\int_0^tL_H(v)\,dv.
\]
Theorem~\ref{pup:low-survival} bounds the right side with an exponent
strictly below $\lambda$.
\end{proof}

\subsection{Removing arrivals after a high-degree crossing}

For an initially low-degree parent, define the hybrid process to agree
with the actual process up to $\tau_M$, and thereafter delete all
parent-to-tag arrivals.  Outward tag selections and all other exterior
transitions may continue; the retained component after $\tau_M$ follows
the inert law.  For an initially high real parent, delete its incoming
channel from time zero.  For an initially inert parent stub, set
$\tau_M=0$ and use inert evolution from the start.  Let $B_M^0(t)$ be the supremum of the hybrid
retained means over all newborn environments.

\begin{theorem}[A subcritical crossing truncation]\label{pup:hybrid}
Suppose $g=\lambda-\mu>0$.  If $M\ge2$ is fixed and
$\delta=1/(M+1)<g$, then
\begin{equation}\label{pup:hybrid-gap}
 \limsup_{t\to\infty}\frac1t\log B_M^0(t)<\lambda.
\end{equation}
\end{theorem}
\begin{proof}
An initially high or inert parent gives the inert newborn mean, with exponent
$\mu$.  For an initially low parent, histories that have not crossed by
$t$ are covered by Theorem~\ref{pup:low-survival}.  Choose
$\gamma<\lambda$ and $K<\infty$ such that $L_M(r)\le Ke^{\gamma r}$.

Split the crossing histories at $\tau_M=\varepsilon t$, where
$\varepsilon>0$ will be small.  The expected core count carried by
crossings in $dr$ is at most $M L_M(r)\,dr$: the pre-crossing parent
degree-gain intensity is at most $M$, and the crossing leaf is added
outside the retained component.  After the crossing the conditional
mean is at most $C(\tau_M)H(t-\tau_M)$.  Thus the late-crossing part
is at most
\[
 M\int_{\varepsilon t}^t L_M(r)H(t-r)\,dr.
\]
For every $\alpha>\lambda$, $H(v)\le K_\alpha e^{\alpha v}$.
The last display has upper exponent at most
$\alpha-(\alpha-\gamma)\varepsilon$, which is below $\lambda$ when
$\alpha$ is sufficiently close to $\lambda$.

For early crossings, use the same finite reference completion and
polynomial forward likelihood reduction as before.  Change parent
degree gains to $h(t)$ only until $\tau_M$; thereafter restore all
genuine transition rates except the deleted parent-to-tag channel.
The retained-and-crossing-time marginal is autonomous: it is the
previous reference law before crossing, then inert evolution.
Its full reference population has birth rate at most $N+1$.

Relative to the fixed genuine finite-tree law, the reference-path
likelihood includes an additional survival factor for the deleted
channel after crossing.  Since the parent's degree is then at least
$M+1$, its missing intensity is at most $\delta$.  On
$\tau_M\le\varepsilon t$, the inverse likelihood therefore obeys
\[
 Z_t^{-1}\le e^{I_M+\delta t}N(\varepsilon t)^M.
\]
With the same $a>0$ and $\theta=a/(a+2)$ as in the preceding proof,
\[
 \mathbb E_{\rm ref}
 [Z_t^{-a};\tau_M\le\varepsilon t]
 \le K e^{a(\delta+M\varepsilon)t}.
\]
Here the density is for the actual law restricted to having no arrivals
in the deleted channel after crossing; its weighted core expectation
is consequently bounded by the unrestricted genuine finite-tree core
mean.  The same three-factor H\"older argument shows that the early
crossing contribution has upper exponent at most
\[
 \theta\mu+(1-\theta)\alpha+
       \theta(\delta+M\varepsilon).
\]
Choose $\varepsilon>0$ with $\delta+M\varepsilon<g$, and then
$\alpha>\lambda$ sufficiently close to $\lambda$ for this exponent
and the late-crossing exponent both to be below $\lambda$.
All forward likelihood factors are polynomial in $t$ and there are
only finitely many starting parent degrees.  This proves the claim.
\end{proof}


\section{A summable disagreement exploration on the evolving old core}
\label{erc:section}

We couple two retained components whose roots differ by one leaf. The proof charges the future families born from affected cores and sums the resulting disagreements over the evolving core tree.

Write
\[
 W_T(s)=\sup_x\mathbb E_x\int_0^T e^{-st}C_{\rm tag}(t)\,dt,
 \qquad
 V_{s,T}(y)=\int_0^T e^{-st}P_t^\infty C(y)\,dt.
\]
The first supremum is over the admissible newborn-environment
class, including finite inert-boundary profiles.  A future child in a
finite component with an inert ancestor belongs to this class.
For finite $T$, core-count Yule domination gives
$W_T(s)\le\int_0^T e^{(1-s)t}\,dt<\infty$.
No finiteness of $W(s)=\sup_TW_T(s)$ will be assumed in the comparison.

There is no enlargement of the supremum hidden in the use of a family
inside a finite inert tree.  Replace its ancestor's inert stub by an
edge to a genuine parent of degree $R$, completed by $R-1$ leaves.
The full degree at the old stub endpoint is unchanged.  Couple all
interior transitions, treating outward choices as null in the inert
copy.  The interiors agree until a parent-to-interior arrival, whose
rate is at most $1/R$ because parent degrees never decrease.  Thus the
probability of disagreement through a fixed horizon $T$ is at most
$T/R$.

The tagged newborn core counts have Yule second moments bounded
uniformly in $R$.  Cauchy--Schwarz therefore bounds the discrepancy in
their expected finite-horizon discounted rewards by a constant
depending only on $s,T$ times $\sqrt{T/R}$, which tends to zero.
Each genuine finite exterior has reward at most $W_T(s)$, so the same
holds in the inert limit.  This argument applies to the fixed finite
state encountered at any charged birth.  It uses only an inert
boundary: it does not identify a constant-rate controlled environment
with an original finite exterior.
If a finite exterior profile has several inert stubs, complete each
one in this way.  For $b$ stubs the disagreement probability is at most
$bT/R$, and the same fixed-profile limit applies.

\subsection{The coupling uses aggregate transitions}

For a core $u$, let $D_u$ be its full degree, $L_u$ its leaf count,
and $c_u^\partial$ its number of inert stubs.  Here there is one inert
stub at the retained root and no others.  The inert core generator has
the following two transition types at $u$:
\begin{align}
 a_u&=L_u+\sum_{v\sim u,\ v\text{ core}}\frac1{D_v},
 &&\text{add one leaf at }u,\label{erc:gain-rate}\\
 b_u&=\frac{L_u}{D_u}\le1,
 &&\text{promote a leaf of }u\text{ to a new core with one leaf}.
 \label{erc:bud-rate}
\end{align}
The second transition changes $L_u$ and the number of core neighbours
by opposite amounts and leaves $D_u$ unchanged.

Couple two copies by maximal coupling of the aggregate leaf-gain
channels at each paired core, and, separately, by maximal coupling of
their promotion channels.  These are distinct continuous-time channels.
A promotion at $u$ and a leaf gain at a neighbour therefore do not
occur simultaneously in this coupling.  This construction has precisely
the required generator in each marginal.  It does not impose a common
source-clock ring between the copies; the source identity is absent
from the core-state transition.  Leaves attached to a given core are
exchangeable, so paired promotions may always be matched as new paired
core children.

Start the copies from $y$ and $y^+$, where $y^+$ has one additional
leaf at the retained root.  Mark that root infected at time zero.
The word ``infected'' is solely a coupling mark.  It never changes
either marginal process.

Orient the core tree away from the retained root.  Common cores born
before their parent's first infection remain eligible for exploration;
all cores present initially are eligible.
When an eligible uninfected core first receives an unmatched leaf gain,
mark it infected permanently.  At its infection, freeze its list of
currently existing common core neighbours.  Every paired or unpaired
new core child born from it \emph{after} its infection is instead
charged as an entire newborn family, permanently excluded from the
exploration.  The actual evolution and feedback of that family remain
present in the coupling.
Every later descendant of a charged family remains in that same
charged family and is also excluded from the exploration.

An eligible uninfected core has identical local state in both copies.
Its only possible infected neighbour is its predecessor toward the
root.  Indeed, the infected part of the eligible common tree is
connected, and a charged family attaches to that tree only at its
already infected parent.  Before an uninfected core is reached, its
common promotions can grow arbitrary common fringes; those fringes
are included when their attachment core is eventually infected.
Thus no initial or subsequently grown common fringe is discarded.

\subsection{A local reservoir controls all return influences}

Fix an infected eligible core $u$ and reset its age to zero at infection.
Let $d\ge2$ be its smaller full degree then.  The two degrees differ
by exactly one at that time.  Define, at subsequent age $t$,
\[
 Z_u(t)=|D_u^+(t)-D_u^-(t)|,\qquad
 R_u(t)=\#\{\text{union promotions at }u\text{ after infection}\},
\]
and let $B_u(t)$ count its frozen outward neighbours subsequently
infected from $u$.  A simultaneous promotion in the two marginals
counts once in $R_u$, even though two newborn families are charged.
Initially $(Z_u,R_u,B_u)=(1,0,0)$.

There are at most $d-1$ frozen outward candidates.  For the initial
root this follows from its inert stub; for all later infected cores
one neighbour is their infected predecessor.  This is the reason the
estimate below is a root-boundary increment estimate.

Put $h(D)=1-1/D$.  Rewriting \eqref{erc:gain-rate} gives
\[
 a_u=D_u-c_u^\partial-
       \sum_{v\sim u,\ v\text{ core}} h(D_v).
\]
Every still-uninfected frozen neighbour makes equal contributions in
the two copies.  The infected predecessor and $B_u$ infected outward
neighbours make at most $(1+B_u)/2$ difference, because
$h(D)\in[1/2,1)$ for a core.  For each union promotion after infection,
a paired child contributes at most $1/2$ difference; a single child
contributes $h(D)<1$.  Consequently, at every time,
\begin{equation}\label{erc:local-discrepancy}
 |a_u^+-a_u^-|
 \le Z_u+R_u+\frac{1+B_u}{2}.
\end{equation}
This bound includes arbitrary repeated feedback through the predecessor
and through all already infected or charged branches.

An unpaired leaf gain changes $Z_u$ by one in absolute value; its total
rate is $|a_u^+-a_u^-|$.  A union promotion increments $R_u$ by one
at rate $\max(b_u^+,b_u^-)\le1$.  For an uninfected frozen outward
neighbour $v$, the only unequal term in its aggregate leaf-gain rate
is the contribution from $u$.  Its infection rate is therefore
\begin{equation}\label{erc:edge-rate}
 \left|\frac1{D_u^+}-\frac1{D_u^-}\right|
 \le \frac{Z_u}{d(d+Z_u)}.
\end{equation}
Full degrees never decrease, which proves the inequality even if the
sign of the degree difference changes.  In particular the birth rate
of $B_u$ is at most $(d-1)Z_u/[d(d+Z_u)]\le1$.

It follows that
\[
 Q_u=Z_u+R_u+B_u
\]
has unit upward jumps with total predictable rate at most
\begin{equation}\label{erc:reservoir-rate}
 Z_u+R_u+\frac{1+B_u}{2}+1+1
 \le Q_u+\frac52.
\end{equation}
Downward jumps can only improve the comparison.  The absence of a
simultaneous promotion and neighbour infection is important here:
it follows from the aggregate-channel coupling, not from coupling
the original source clocks.

Conditionally at infection, $Q_u$ is thus dominated by a pure birth
process $Y$ started at $1$, with transition rate
\begin{equation}\label{erc:dominating-chain}
 n\longrightarrow n+1\quad\text{at rate }n+\frac52.
\end{equation}
For example, a common Poisson measure constructs this domination,
since the upward rate bound is increasing in the current state;
downward jumps of $Q_u$ preserve the order.  The comparison is uniform
in the whole prior history and all states of the attached branches.

\subsection{An exact rational reproduction certificate}

For $d\ge2$ and $s>0$ define
\begin{equation}\label{erc:mass}
 m_d(s)=\frac{d-1}{d}\,
 \mathbb E\int_0^\infty e^{-st}\frac{Y(t)}{d+Y(t)}\,dt,
 \qquad m(s)=\sup_{d\ge2}m_d(s).
\end{equation}
The compensator of the infections from $u$, together with
\eqref{erc:edge-rate} and the increasing function $z/(d+z)$, gives
\begin{equation}\label{erc:offspring}
 \mathbb E\left[
   \sum_{v:\,u\to v}e^{-s(\tau_v-\tau_u)}
      \,\middle|\,\mathcal F_{\tau_u}\right]
 \le m_d(s)\le m(s).
\end{equation}
No independence of different infected cores is asserted or needed.

\begin{lemma}[A rational discount below the original exponent]
\label{erc:certificate}
For every $s\ge s_0:=11/20$,
\begin{equation}\label{erc:subcritical-mass}
 m(s)\le\frac{19}{20}<1.
\end{equation}
\end{lemma}
\begin{proof}
It suffices to use $s=s_0$, by monotonicity in the discount.
Let $w_n$ be the expected discounted occupation time of state $n$ for
\eqref{erc:dominating-chain}.  The exponential holding times give
\[
 w_1=\frac{20}{81},\qquad
 w_{n+1}=w_n\frac{20n+50}{20n+81},\qquad
 \sum_{n\ge1}w_n=\frac{20}{11}.
\]
Consequently, for every integer $N\ge1$,
\begin{equation}\label{erc:rational-upper}
 m_d(s_0)
 =\frac{d-1}{d}\left(\frac{20}{11}
                  -d\sum_{n\ge1}\frac{w_n}{d+n}\right)
 \le\frac{d-1}{d}\left(\frac{20}{11}
                  -d\sum_{n=1}^N\frac{w_n}{d+n}\right).
\end{equation}
With $N=16$, exact rational comparison for $2\le d\le63$ gives a
maximum of the displayed upper bounds at $d=6$, equal to
\[
 \frac{1669075573443141901662377825}
      {1772128093509182310390306279}
 <\frac{19}{20}.
\]
The complete finite calculation, using only integer arithmetic and
rational fractions, is in \texttt{scripts/verify\_infection\_kernel.py}.

For the remaining degrees, the mean of $Y$ is
$\mathbb EY(t)=(7/2)e^t-5/2\le(7/2)e^t$.  Hence, for $d\ge64$,
\begin{align*}
 m_d(s_0)
 &\le\int_0^\infty e^{-t/2}
            \min\{1,(7/(2d))e^t\}\,dt\\
 &=4\sqrt{\frac7{2d}}-\frac7d
 \le\sqrt{\frac78}<\frac{15}{16}<\frac{19}{20}.
\end{align*}
The first inequality uses $s_0\ge1/2$ and
$Y/(d+Y)\le\min\{Y/d,1\}$; the expectation of this minimum is
bounded by the minimum of the expectation and one.  The last radical
comparison is equivalently the rational inequality
$7/8<(15/16)^2$.  This covers every integer degree.
\end{proof}

Generation by generation in the infection tree, conditional expectation
and \eqref{erc:offspring} now give
\begin{equation}\label{erc:total-infections}
 \mathbb E\sum_{u\text{ infected}}e^{-s\tau_u}
 \le\sum_{k\ge0}m(s)^k\le20
 \qquad(s\ge s_0).
\end{equation}
For a fully formal construction, first stop at finite population and
infection generation, use the compensator identities there, and then
pass to the limit by Tonelli and monotone convergence.  All finite
horizons are nonexplosive by real-population Yule domination.

\subsection{The old-state boundary gradient}

\begin{theorem}[One boundary leaf has a uniform resolvent effect]
\label{erc:gradient}
For every finite retained inert state $y$, every finite $T$, and
$s\ge11/20$, addition of one leaf at its root satisfies
\begin{equation}\label{erc:gradient-bound}
 \boxed{
 |V_{s,T}(y^+)-V_{s,T}(y)|
 \le\frac{2W_T(s)}{s(1-m(s))}
 \le\frac{40}{s}W_T(s).
 }
\end{equation}
The bound is uniform in the number of old cores, their degrees, and
their attached finite fringes.
\end{theorem}
\begin{proof}
Every eligible common core is present in both copies for all time
after its paired birth.  Those cores contribute equally to the two
core counts, whether or not they are eventually infected.  Every other
core belongs to a unique charged family: take its first ancestor born
from an eligible infected core after that core's infection.  These
families are disjoint within each marginal.  This partitions all
possible contributions to the absolute core-count difference.

At each union promotion from an infected core, at most two newborn
families are charged.  Conditionally on their birth time and the joint
present, each marginal has the correct future generator, so each
remaining discounted reward is bounded by $W_T(s)$.  This assertion
does not require independence from the other family or from the
infection exploration.  The birth time is a stopping time and every
such family starts in an admissible newborn environment.

The union-promotion rate at $u$ is at most one.  Its expected discounted
charge is therefore at most
\[
 2W_T(s)\,\mathbb E\int_{\tau_u}^T e^{-st}\,dt
 \le\frac{2W_T(s)}s\,\mathbb E e^{-s\tau_u}.
\]
Sum over infected eligible cores and use
\eqref{erc:total-infections}, or its sharper bound $1/(1-m(s))$.
Pathwise domination of the absolute count difference by the charged
family counts and Tonelli's theorem give \eqref{erc:gradient-bound}.
Every occurrence of $W_T$ has finite horizon, so the argument is not
assuming the uniform resolvent conclusion it will prove.
\end{proof}

\section{Closing the uniform-growth gap}
\label{erc:closure-section}

For an actual newborn environment, let $\tau_M$ be the first time its
external parent degree exceeds $M$, taking $\tau_M=0$ if it is already
larger or if the parent is an inert stub.  Let the hybrid agree with the actual process through
$\tau_M$, and thereafter delete parent-to-root arrivals.  Write
$A_M(t)$ for the supremum of its retained core mean.  Theorem~\ref{pup:hybrid} proves that if
$\mu<\lambda$ and $\delta=1/(M+1)<\lambda-\mu$, then
\begin{equation}\label{erc:hybrid-assumption}
 \gamma_M:=\limsup_{t\to\infty}t^{-1}\log A_M(t)<\lambda.
\end{equation}
This theorem includes arbitrary old states at the
crossing.  Only this strict gap for each sufficiently large fixed $M$
is needed below; no bound on its size uniform in $M$ is required.

\begin{proposition}[A return coefficient independent of the exit mass]
\label{erc:absorb}
For every $s\ge11/20$ and finite $T$,
\begin{equation}\label{erc:absorb-finite}
 W_T(s)\le\int_0^T e^{-st}A_M(t)\,dt
                 +\frac{40}{s^2(M+1)}W_T(s).
\end{equation}
In particular, if $s>\gamma_M$ and $40/[s^2(M+1)]<1$, then
\begin{equation}\label{erc:absorb-bound}
 W(s)\le
 \frac{\int_0^\infty e^{-st}A_M(t)\,dt}
      {1-40/[s^2(M+1)]}<\infty,
 \qquad\lambda\le s.
\end{equation}
\end{proposition}
\begin{proof}
After the crossing, the actual retained component has conditional
generator $\mathcal L_\infty+q_t\Delta_{\rm root}$, where
$0\le q_t\le\delta$.  For $f_r(y)=V_{s,r}(y)$, the finite-horizon
resolvent identity is
\[
 \partial_rf_r=C+(\mathcal L_\infty-s)f_r,\qquad f_0=0.
\]
Apply the time-dependent discounted Dynkin identity to
$e^{-st}f_{T-t}(Y_t)$ from $\tau_M\wedge T$ through $T$, with no
post-crossing term on $\{\tau_M\ge T\}$.  The reward before
$\tau_M$ and the inert continuation value at $\tau_M$ give exactly
the hybrid reward through $T$.  The remaining term is
\[
 \mathbb E_x\int_{\tau_M}^T e^{-st}q_t
             \Delta_{\rm root}V_{s,T-t}(Y_t)\,dt.
\]
The theorem above, $W_{T-t}(s)\le W_T(s)$, and
$\int_0^T e^{-st}dt\le1/s$ bound this term from above by
$40\delta W_T(s)/s^2$.  Localization is justified by finite-horizon
population moments; the resulting boundary increment is already
uniform in the old state.  Taking the supremum over initial newborn
environments proves \eqref{erc:absorb-finite} without interchanging a
supremum and an integral.

The coefficient is strictly less than one, so division is legitimate
at finite $T$.  Monotone convergence then proves finiteness of $W$.
Finally each $F_x(t)$ is nondecreasing, and hence
$F_x(t)\le se^{st}\int_0^\infty e^{-su}F_x(u)du\le se^{st}W(s)$.
Taking the supremum over $x$ gives $\lambda\le s$.
\end{proof}

\begin{theorem}[Equality of the original and uniform exponents]
\label{erc:equality}
The original friend-tree mean exponent and the uniform newborn mean
exponent agree:
\begin{equation}\label{erc:lambda-mu}
 \boxed{\lambda=\mu.}
\end{equation}
\end{theorem}
\begin{proof}
The existing comparison gives $\mu\le\lambda$, and the certified lower
bound on the original exponent gives $\mu\ge0.565>11/20$.
Suppose $\mu<\lambda$.  Choose a fixed integer $M\ge132$ large enough
that $1/(M+1)<\lambda-\mu$.  Then \eqref{erc:hybrid-assumption} holds.
Choose
\[
 \max\{11/20,\gamma_M\}<s<\lambda.
\]
The elementary rational inequality
$40/[(11/20)^2\cdot133]=16000/16093<1$ ensures
$40/[s^2(M+1)]<1$.  Proposition~\ref{erc:absorb} therefore yields
$\lambda\le s$, contradicting $s<\lambda$.
\end{proof}


\section{Effective computation of the exact exponent}
\label{ecc:section}

\subsection{The two exact characterizations}

For a newborn core in an admissible finite environment $x$, let
\[
 F_x(t)=\mathbb E_x C_{\rm tag}(t),\qquad
 G(t)=\sup_x F_x(t),\qquad H(t)=1+\int_0^tG(u)\,du.
\]
The supremum includes the finite inert environments used by the pruning
construction; their finite-horizon rewards are limits of rewards in
genuine finite trees.  The future-child decomposition in Section~\ref{uniform:section} gives
$H(t+u)\le H(t)H(u)$ and $G(t)\le H(t)\le1+tG(t)$.
Theorem~\ref{erc:equality} therefore gives
\begin{equation}\label{ecc:inf-characterization}
 \boxed{\displaystyle
 \mu=\inf_{T\in\mathbb Q_{>0}}\frac{\log H(T)}{T}.}
\end{equation}
Restricting to rational times is valid by continuity of $H$.
Every fixed rational-time value $H(T)$ is effectively computable.

There is also the exact scalar pressure equation from the inert newborn
process.  Give its initial root core genealogical depth zero.  Let
$\sigma_v$ be the birth time of a future core $v$, and put
\[
 A_m^\infty(s)=\mathbb E_\infty
       \sum_{v:\,\operatorname{depth}(v)=m}e^{-s\sigma_v},
 \qquad m\ge1,\quad s>0.
\]
Theorem~\ref{gc:root} proves that
\begin{equation}\label{ecc:pressure}
 r_\infty(s)=\lim_{m\to\infty}A_m^\infty(s)^{1/m}
\end{equation}
exists, is positive, continuous, and strictly decreasing, and that
\begin{equation}\label{ecc:unique-root}
 \boxed{r_\infty(\mu)=1.}
\end{equation}
Moreover,
\begin{equation}\label{ecc:lower-witness}
 s<\mu\quad\Longleftrightarrow\quad
 A_m^\infty(s)>12(s+3)\text{ for some finite }m.
\end{equation}
The older constant $12(s+3)$ is sufficient here; no sharper minorization
constant is needed for computability.

The two characterizations supply complementary strict certificates. No equality of distinct pressure functions away from their critical points is required.

\subsection{What the finite upper calculation consists of}

Here are the effective ingredients for $H(T)$, summarized from Section~\ref{profile:time} so
that the infimum in~\eqref{ecc:inf-characterization} is usable.
Prune the newborn exterior at depth $R$ or its first vertex of degree
at least $M$.  The resulting profile set is finite and explicitly
enumerable.  For all $0\le t\le T$, the locality coupling bounds
the error between the true supremum and the maximum of the retained
profile means by
\begin{equation}\label{ecc:profile-error}
 e^T\left(\frac{T^R}{R!}+\frac{e^T-1}{M}\right).
\end{equation}
For a fixed profile with $n$ initial real vertices, stop above total
population $N$.  The remaining finite generator has rational rates.
Its matrix exponential is computed by uniformization, using a rational
Poisson-series enclosure.  The omitted mean is at most
\begin{equation}\label{ecc:population-error}
 \frac{n(n+1)e^{2T}-ne^T}{N}.
\end{equation}
Finally $G$ is $e^T$-Lipschitz on $[0,T]$, so rational-grid quadrature,
together with~\eqref{ecc:profile-error}--\eqref{ecc:population-error},
computes $H(T)$ with any specified positive rational error.
All exponential and logarithmic quantities in these bounds admit
elementary certified rational enclosures.

Consequently the following upper search is sound and complete:
enumerate positive rational $T$ and improving certified bounds
$H(T)\le\overline h$, and stop when
\begin{equation}\label{ecc:upper-witness}
 \overline h<e^{sT}
\end{equation}
has been certified by rational enclosures.  Success implies $\mu<s$.
For every rational $s>\mu$, submultiplicativity and
\eqref{ecc:inf-characterization} supply a strict witness, so this search
eventually stops.  It is not claimed to stop when $s=\mu$.

For the lower search, enumerate finite histories of the exactly
eliminated waiting-event kernel.  At rational $s$, all retained beta
weights are rational and nonnegative.  The earlier program
\texttt{scripts/inert\_growth\_kernel.py} computes increasing finite sums
$L_{m,R,N}(s)\le A_m^\infty(s)$.  Enumerating $m,R,N$ and stopping when
\[
 L_{m,R,N}(s)>12(s+3)
\]
is therefore sound and, by~\eqref{ecc:lower-witness}, complete for
$s<\mu$.  Each attempted certificate is finite exact arithmetic.

\begin{theorem}[Effective two-sided computation]\label{ecc:computable}
For every positive rational $\varepsilon$, there is an algorithm which
halts and outputs rational $a<b$ with
\[
 a<\mu<b,\qquad b-a<\varepsilon.
\]
Thus the original random friend-tree exponent is a computable real.
\end{theorem}
\begin{proof}
Enumerate pairs of positive rationals $a<b$ with $b-a<\varepsilon$.
Dovetail the lower search at $a$ and the upper search at $b$ for every
pair: at each finite stage perform finitely many additional steps of
finitely many searches, ensuring that each fixed search receives
unboundedly many steps.  Return the first pair for which both
certificates have been found.

Each returned interval is sound.  Density of the rationals supplies
a pair $a<\mu<b$ of the required width.  For that pair both searches
halt, so the dovetailed algorithm halts.  The algorithm never needs to
decide whether a proposed rational test value equals $\mu$.
\end{proof}

\subsection{A convergent rational series with a certified remainder}

Start from $a_0=1/2$ and $b_0=3/5$, which strictly bracket the known
exponent.  At stage $n\ge1$, use the preceding search restricted to
\[
 a_{n-1}<a_n<b_n<b_{n-1},\qquad b_n-a_n<2^{-n}.
\]
Because the previous interval strictly contains $\mu$, every stage
terminates.  Fix a computable ordering and dovetailing schedule, so
the rational sequence is uniquely specified by this algorithm.
It follows that
\begin{equation}\label{ecc:rational-series}
 \boxed{\displaystyle
 \mu=\frac12+\sum_{n=1}^{\infty}(a_n-a_{n-1}),\qquad
 0<\mu-a_n<2^{-n}.}
\end{equation}
All summands are positive rationals obtained by finite certificate
searches.  The midpoint $q_n=(a_n+b_n)/2$ satisfies
$|q_n-\mu|<2^{-n-1}$.

This is an algorithmic convergent series with an explicit remainder
bound after each completed stage.  It is not a claimed closed formula
for the summands or an efficient high-precision numerical method.
The search may be extremely expensive, and this work gives no useful
running-time estimate as $\varepsilon\downarrow0$.  Those qualifications
do not weaken the termination or error guarantees.


\section{The canonical resolvent and its exact kernel series}
\label{canonical:section}

Theorem~\ref{erc:equality} gives $W(s)<\infty$ for every $s>\mu$.
Pass to the limit $T\to\infty$ in Theorem~\ref{erc:gradient} to obtain
\begin{equation}\label{erc:infinite-gradient}
 |V_s(y^+)-V_s(y)|\le\frac{40}{s}W(s),\qquad s>\mu.
\end{equation}
A certified block $e^{-sT}H(T)\le q<1$ supplies the explicit bound
\begin{equation}\label{canonical:W-bound}
 W(s)\le\int_0^\infty e^{-st}H(t)\,dt
 \le\frac1{1-q}\int_0^T e^{(1-s)t}\,dt.
\end{equation}
Such a block can be found by the complete search after
Theorem~\ref{sc:window}, now with $\lambda=\mu$.

The old-core decomposition gives
\begin{equation}\label{canonical:linear}
 \frac{C(y)}s\le V_s(y)
 \le\frac{C(y)}s(1+W(s)).
\end{equation}
For the upper inequality, partition the future population by the
currently existing cores and their future newborn families, whose
birth intensities are at most one. Integrating their discounted
rewards gives the right side. The lower inequality uses monotonicity
of the core count. These estimates, finite-horizon population moments,
and localization justify the resolvent generator equation
\begin{equation}\label{canonical:generator}
 (s-\mathcal L_\infty)V_s=C.
\end{equation}

The block also controls the remaining time integral effectively:
\begin{equation}\label{canonical:time-tail}
 \int_{nT}^\infty e^{-st}P_t^\infty C(y)\,dt
 \le\frac{C(y)q^n}{1-q}\int_0^T e^{(1-s)r}\,dr.
\end{equation}
Together with a finite-state calculation on $[0,nT]$, this computes
$V_s(y)$ at every finite state, once a block with $e^{-sT}H(T)<1$ has been certified.

\begin{proposition}[Exact eliminated-kernel representation]
\label{canonical:series}
For $s>\mu$, let $Q_s$ be the exact core-ring kernel in
\eqref{model:Q}, and let
\[
 g_s(y)=\frac{C(y)}{C(y)+s}.
\]
Then
\begin{equation}\label{canonical:kernel-series}
 V_s(y)=\sum_{j=0}^\infty(Q_s^jg_s)(y).
\end{equation}
This is a nonnegative convergent series with explicitly given beta
kernel entries. Its value is effectively computable at every finite
state and rational $s>\mu$. In particular $V_s$ is the canonical
positive supersolution for the eliminated kernel, with the uniform
root-leaf increment \eqref{erc:infinite-gradient}.
\end{proposition}
\begin{proof}
Let $\tau_j$ be the successive core-ring times, with $\tau_0=0$.
Between core rings $C$ is constant. The expected discounted reward
before the first ring is $R_sC=g_s$, and the discounted law of the
post-ring state is $Q_s$. The Markov property identifies the expected
reward over $[\tau_j,\tau_{j+1})$ with $Q_s^jg_s$. The original
root's rate-one clock ensures infinitely many rings over an infinite
time interval; nonexplosion ensures $\tau_j\to\infty$. Tonelli's
theorem partitions the total reward into these intervals and proves
\eqref{canonical:kernel-series}. Finiteness follows from
\eqref{canonical:linear}; effective evaluation follows from
\eqref{canonical:time-tail} and finite-state uniformization. The
first-step equation is $V_s=g_s+Q_sV_s$, and
\eqref{erc:infinite-gradient} gives the boundary increment.
\end{proof}

The finite-window construction also gives a computable supersolution
with bounded increments at \emph{every} ordinary birth, rather than
only at a root-leaf addition. At every rational $s>\mu$, a complete
search finds $T,q$ with $e^{-sT}H(T)\le q<1$, and
Theorem~\ref{sc:window} gives $U_{s,T}$ and its explicit increment
constant. The canonical and finite-window constructions are both
useful: the former solves the exact resolvent equation, while the
latter has a direct finite-time formula and controls all ordinary
birth increments.

\section{Every bounded predictable input has the same mean exponent}
\label{bcl:section}

Fix $Q<\infty$, and let $\mathcal Q_Q$ be the nonanticipating policies
that add a leaf at the retained root with predictable intensity
$0\le q_t\le Q$.  Apart from these arrivals the retained component
has the inert generator.  Policies may depend on the observed component
and additional nonanticipating randomness.  They do not anticipate
the future Poisson driving measures.  The component starts from the
inert newborn state $x_\infty$.  Put
\[
 F_q(t)=\mathbb E_q C(t),\qquad
 F_Q^-(t)=\inf_{q\in\mathcal Q_Q}F_q(t),\qquad
 F_Q^+(t)=\sup_{q\in\mathcal Q_Q}F_q(t).
\]
The resolvent estimate controls $F_Q^+$; a preparation and isolation argument gives the matching lower rate uniformly over the same class. No ordering of the tree laws under different input rates is required.

\begin{lemma}[Uniform preparation and isolation of an inert descendant]
\label{bcl:preparation}
For every integer $M\ge2$, every $q\in\mathcal Q_Q$, and every $t\ge0$,
\begin{equation}\label{bcl:uniform-lower}
 \boxed{
 F_q(t+2)\ge p_{M,Q}e^{-t/M}F_\infty(t),\qquad
 p_{M,Q}=\frac{M-1}{M}\,e^{-2Q-5}
                         (1-e^{-1})^{M-2}>0.
 }
\end{equation}
Here $F_\infty$ is the inert newborn mean, already known to have
exponent $\mu$.  Only a fixed preparation interval suppresses the
external input.  No suppression of that input is imposed after time
two.
\end{lemma}
\begin{proof}
Write $u$ for the original retained root.  Realize the controlled
arrivals by thinning an independent rate-$Q$ candidate process,
accepting a candidate with conditional probability $q_t/Q$ when
$Q>0$.  For $Q=0$ this channel is absent.  Suppress all candidates
on $[0,2]$, at probability cost $e^{-2Q}$, independently of the
internal driving clocks.  On this event the preparation is the same
for every admissible policy.

In the first unit interval keep $u$'s clock silent and require its
initial one-leaf Yule family to have exactly $M-1$ leaves at time one.
The root silence has probability $e^{-1}$, and the Yule probability is
\[
 e^{-1}(1-e^{-1})^{M-2}.
\]
The full degree of $u$ at time one is then $M$, counting its inert
stub.  There are no other cores in this prepared component.

In the second unit interval require exactly one ring of $u$'s clock.
This has probability $e^{-1}$.  Before that ring all its real
neighbours are leaves, whose number can only increase, so the
probability of choosing a leaf is at least $(M-1)/M$.  This promotion
creates a core child $v$ with one leaf.  Require $v$ and its new leaf
to have no rings before time two.  Given the promotion time, this has
probability at least $e^{-2}$.  Thus at time two $v$ is exactly a
newborn retained core, attached to the real parent $u$ of degree at
least $M$.  Multiplication of these conditional probability bounds
gives the displayed $p_{M,Q}$.  Other leaves at $u$ are free to grow
throughout the second interval; this cannot decrease its full degree.

Condition on the whole preparation history at time two.  Split $u$'s
future rate-one clock into a rate-$1/M$ candidate stream for the edge
$u\to v$ and a rate-$(1-1/M)$ residual stream, with independent fresh
increments.  At a candidate ring choose $v$ with probability
$M/D_u$.  If the candidate is rejected, choose uniformly among the
other $D_u-1$ neighbours; at a residual ring use the same uniform
choice among those other neighbours.  This is a correct construction
of the original source choices: the rate to $v$ is $1/D_u$, and the
rate to any other neighbour is
\[
 \frac{(1/M)(1-M/D_u)+(1-1/M)}{D_u-1}
 =\frac1{D_u}.
\]
It remains valid forever because degrees are nondecreasing and
$D_u\ge M$.

The probability that this candidate stream has no ring in the next
$t$ units of time is exactly $e^{-t/M}$.  The stream is independent
of the internal drivers of $v$ and its future descendants.  On its
no-ring event there can be no arrival from $u$ into $v$.  The retained
descendant component at $v$ consequently follows exactly the inert
newborn process.  An outward choice $v\to u$ changes only the parent
and is a null transition for this descendant component.  The external
control, even if it reacts to the descendant history, also changes
only $u$ and cannot send an arrival across the edge without a point
of the reserved candidate stream.

Conditional on the preparation history, the expected core reward of
this descendant component multiplied by the no-candidate indicator is
therefore $e^{-t/M}F_\infty(t)$.  Its cores are all counted in the
original retained component.  Averaging over the preparation event
proves~\eqref{bcl:uniform-lower}.
\end{proof}

\begin{theorem}[Matching upper and lower rates uniformly over controls]
\label{bcl:matching}
For every fixed finite $Q\ge0$,
\begin{equation}\label{bcl:inf-sup}
 \boxed{
 \lim_{t\to\infty}\frac{\log F_Q^-(t)}t
 =\lim_{t\to\infty}\frac{\log F_Q^+(t)}t
 =\mu.
 }
\end{equation}
In particular, every single admissible predictable policy has mean
exponent $\mu$.  The conclusion also holds for horizon-dependent
choices of policies with the same common bound $Q$.
\end{theorem}
\begin{proof}
For fixed $M$, Lemma~\ref{bcl:preparation} and the inert exponent give
\[
 \liminf_{t\to\infty}\frac{\log F_Q^-(t)}t
 \ge\mu-\frac1M.
\]
The fixed time shift and the positive preparation probability disappear
at the exponential scale.  Let $M\to\infty$.

For completeness, the upper bound uses the canonical resolvent
$V_s$, not a monotonicity claim.  At any $s>\mu$, write
$K_s=40W(s)/s$ and
\[
 U_s=V_s+\frac{QK_s}{s}.
\]
The established root increment bound and resolvent identity give
\[
 (s-\mathcal L_\infty-q\Delta_{\rm root})U_s\ge C
 \qquad(0\le q\le Q),\qquad U_s\ge C/s.
\]
Stopped Dynkin and nonnegativity give
\[
 F_Q^+(t)\le sU_s(x_\infty)e^{st}.
\]
The controlled population is nonexplosive, since its birth rate is at
most its real population plus $Q$.  Fatou's lemma removes the
localization.  Let $s\downarrow\mu$.  Since
$F_Q^-\le F_q\le F_Q^+$, both the uniform statement and the assertion
for each policy follow.
\end{proof}

\begin{corollary}[A permanently fixed parent has exponent $\mu$]
\label{bcl:fixed-parent}
If the parent has permanently fixed degree $d\ge2$, supplying its
newborn child with independent rate-$1/d$ arrivals, and $F_d(t)$ denotes
the retained mean, then
\[
 \boxed{\lim_{t\to\infty}\frac1t\log F_d(t)=\mu.}
\]
\end{corollary}
\begin{proof}
Take the constant policy $q_t=1/d$ in
Theorem~\ref{bcl:matching}.
\end{proof}

The argument extends to every fixed finite retained starting component
with one inert parent stub and a core root.  Its upper constant is
$sU_s(y)$.  For the lower bound, a finite positive-probability
preparation first creates a root leaf if necessary, then builds a
large root degree and a fresh child as above.  Such a root has at
least one real neighbour, and prescribing a choice of the root by
that neighbour has positive probability.  Suppressing rate-$Q$
candidates during this finite preparation keeps its probability
uniform in the policy.  The resulting constant may depend on the
fixed initial component; the exponent remains $\mu$.


\section{A refined rational motif certificate}
\label{numerical:section}

This section uses the upper motif certificate from~\cite{UzayTables}.
The numerical data are pinned separately from the later revision used
for the existence proof. The unchanged tables and their original
all-state drift decomposition are inputs; the refinement consists of
a correlated upper weight and a constant correction per core.

\subsection{The potential and its grouped upper weight}

Root the real core tree at a core vertex. A core has type $(L,k,h)$,
where $L$ is its leaf count, $k$ its number of core children, and
$h\in\{0,1\}$ indicates a core parent. The original potential is
\begin{align}\label{num:potential}
 \Phi(T)={}&\sum_v f(v)+\sum_{p\to w}g(p,w)
       +\sum_{p\to u\to w}h_V(p,u,w)\notag\\
       &+\sum_u\sum_{\{w_1,w_2\}\subset\operatorname{ch}_K(u)}
                       h_C(u;w_1,w_2).
\end{align}
All tables have denominator $10^{10}$. The vertex table saturates at
$L_F=800$ and $K_F=150$. For an endpoint of full degree $D\ge24$,
the interpolation measure is
\[
 \pi_{D,k,h}=(1-24/D)\delta_{\mathrm{INF}}
       +(24/D)\delta_{T_1(\min(k,13),h)}.
\]
The vertical middle has at most eight core children and degree below
$16$; the cherry middle has at most six core children and degree below
$12$. Terms outside these middle domains vanish. Endpoint evaluations
are multilinear in the displayed measures. The exact arrays and
projection rules are generated by the pinned \texttt{cert\_build.c}.

Assign an edge and a vertical path to their bottom core, and a cherry
to its middle core. Let $W(v)$ be the resulting contribution at $v$,
so that $\Phi=\sum_vW(v)$. If $\mathcal P,\mathcal W$ are the original
finite parent and child endpoint sets, put
\begin{align*}
 R(w)=\max\bigl\{0,&\ \max_{p\in\mathcal P}g[p,w],\\
 &\max_{u\in H_u,\ p\in\mathcal P}
                 (g[u,w]+h_V[p,u,w])\bigr\},\\
 H^+(u)&=\max\{0,\max_{w_1,w_2\in\mathcal W}h_C[u,w_1,w_2]\}.
\end{align*}
Here $H_u$ is the supported vertical-middle set and $C_u$ below is
the supported cherry-middle set. Define
\begin{equation}\label{num:Q}
 Q(v)=f(v)+h\sum_w\pi_v(w)R(w)
       +\mathbf1_{\{v\in C_u\}}\binom{k}{2}H^+(v).
\end{equation}
The edge and vertical terms in $R$ use the same parent. Positivity of
the interpolation weights therefore gives the pointwise bound
$W(v)\le Q(v)$. Exact evaluation gives
\begin{equation}\label{num:comparison}
 m C\le\Phi\le\overline A C,
 \qquad m=\frac{2989767809}{10^{10}},
 \qquad \overline A=\frac{8424600447}{5000000000}.
\end{equation}
The lower bound is part of the original certificate. The upper bound
uses $\sup_vQ(v)=\overline A$.

Write $\nu=1133/2000$, and let $B_f(v;T)$ be the original assigned
vertex drift. Its vertex reserve is
$\delta(v;T)=\nu f(v)-B_f(v;T)$. The verified nonnegative edge,
vertical, and cherry residuals imply the all-tree inequality
\begin{equation}\label{num:original-drift}
 \mathcal L\Phi(T)\le\nu\Phi(T)-\sum_{v\in K(T)}\delta(v;T).
\end{equation}
This is the drift decomposition to which the refinement is applied.

\subsection{A constant correction per core}

Set
\begin{gather}\label{num:parameters}
 \beta=\frac{2832242711}{5000000000},\qquad
 \tau=\frac{5294970481}{10^{14}},\qquad
 \Phi_\tau=\Phi+\tau C,\\
 \zeta=\frac{199752266575452391}
              {1200500000000000000000000000}>0.
 \label{num:reserve}
\end{gather}

\begin{theorem}[Corrected full-tree drift]\label{num:drift}
Every finite full tree with at least one core satisfies
\begin{equation}\label{num:corrected-drift}
 \mathcal L\Phi_\tau\le\beta\Phi_\tau-\zeta C.
\end{equation}
Consequently $\mu\le\beta=0.5664485422$.
\end{theorem}
\begin{proof}
Put $b_v=L_v/D_v$. Only promotion increases the core count, so
$\mathcal LC=\sum_v b_v$. Since $\nu>\beta$ and $W(v)\le Q(v)$,
\eqref{num:original-drift} gives
\begin{align*}
 \beta\Phi_\tau-\mathcal L\Phi_\tau
 &\ge\sum_v\left[\delta(v;T)-(\nu-\beta)W(v)
                            +\tau(\beta-b_v)\right]\\
 &\ge\sum_v\left[\delta(v;T)-(\nu-\beta)Q(v)
                            +\tau(\beta-b_v)\right].
\end{align*}
The exact additional finite certificate checks
\begin{equation}\label{num:row-check}
 \delta_i-(\nu-\beta)Q_i+\tau(\beta-b_i)\ge\zeta
\end{equation}
in all $484138$ vertex rows. The minimum is the displayed $\zeta$.

For completeness, the original tail interpolation also proves that
these rows cover all degrees. In the unsaturated outer box only the
incoming leaf-gain intensity varies; the expression is affine in
that intensity, and both endpoints are checked. On the leaf tail the
enclosing parameter rectangle is
\[
 b\in\left[\frac{801}{801+k+h},1\right],\qquad
 t\in\left[0,\frac{24}{801+k+h}\right].
\]
The vertex value is constant there, the assigned drift is bilinear in
$(b,t)$, and $Q$ is affine in $t$. The correction is affine in $b$.
Thus \eqref{num:row-check} is bilinear and follows from its four corner
checks. The post-promotion edge uses child-count label $k+1$ and $Q$
uses the pre-promotion label $k$, at the same value of $t$. The exact
$L=800$ boundary is checked separately.

For overflow $k>150$, the rectangle is
$b\in[0,1]$, $t\in[0,24/(151+h)]$, with the clamped child-count
label $13$. The same four-corner argument applies. Hence
\eqref{num:row-check} holds at every actual vertex, and summation
proves \eqref{num:corrected-drift}.

The lower comparison $(m+\tau)C\le\Phi_\tau$ and the upper comparison
$(\overline A+\tau)C\ge\Phi_\tau$ follow from
\eqref{num:comparison}. Dropping the strictly positive reserve and
applying the resulting drift inequality gives expected growth at most
$e^{\beta t}$ from every fixed finite seed. Equation
\eqref{model:imported-exponent} identifies the original exponent,
so $\mu\le\beta$.
\end{proof}

The exact feasible interval for $\tau$ at this fixed $\beta$ is
\[
 \frac{763278221264673581}{14415721236665000000000}
 \le\tau\le
 \frac{32452312724432471239}{612865341021362000000000}.
\]
The lower constraint is the outer row $(L,k,h)=(143,150,1)$ at its
larger incoming-rate endpoint, with $b=143/294$; the upper constraint
is the analogous row $(795,150,1)$, with $b=795/946$. The chosen
rational $\tau$ is strictly inside this interval. No optimality claim
for this pair is needed.

At the deterministic three-vertex seed,
$\Phi(T_3)=14606753031/10^{10}$ and $C(T_3)=1$. Dividing the generator
inequality by the current population and iterating gives the explicit
finite-size bound
\begin{equation}\label{num:finite-size}
 \mathbb E C(T_n)\le
 \frac{14606753031/10^{10}+\tau}{m+\tau}
 \prod_{j=3}^{n-1}\left(1+\frac\beta j\right),\qquad n\ge3.
\end{equation}

\subsection{Reproduction of the certificate}

The repository supplies \texttt{scripts/verify\_bound.py}, which fetches
the immutable upstream inputs, verifies their hashes, builds the
original rational binary, runs the complete GMP verifier, and checks
the corrected rows in exact arithmetic. The upstream manifest and
the two result files are under \texttt{verification/}. The full run
checks $17329108$ original motif inequalities and $37$ finite trees,
then all $484138$ corrected vertex rows. The tail argument above
extends those finite comparisons to arbitrary degrees.

The corrected row check reads the original arrays; it does not replace
the original edge, vertical, or cherry checks. All decisions are
rational sign comparisons. Floating-point calculations may propose
parameters but are not part of the certificate. Reproduction with the
same verified source is a reproducibility check, rather than an
independently implemented second verifier.

\section{Computational scope and further questions}
\label{discussion:section}

The characterization $r_\infty(\mu)=1$ uses the original evolving
environment. Its coefficients can be computed by exact positive
histories, and the time-block upper certificates have explicit
finite-time errors. These two features make $\mu$ a computable real.
They do not imply that the resulting universal certificate search is
an efficient numerical solver. The practical numerical assertion of
this paper is the finite enclosure in Theorem~\ref{main:bound}.

The arithmetic question remains separate. A computable unique root
can be rational, algebraic irrational, or transcendental. Positivity
of the eliminated kernel supplies no arithmetic classification by
itself. To prove irrationality of the original $\mu$, one would need
an arithmetic specialization or approximation theorem for its actual
characteristic structure, with a nonvanishing statement and a
sufficient denominator or height bound. No reduction of the present
infinite evolving-environment pressure to a finite polynomial relation
in one Mellin--exponential function is proved here.

The companion local-model draft~\cite{Local} has an independent
branching representation and an explicit scalar Mellin equation. Its
arithmetic theorem applies to those modified models. It cannot be
transferred to the original friend tree by numerical proximity of
their exponents or similarity of their local transition formulas.

The remaining probabilistic questions include sharp mean asymptotics,
effective and practical high-precision evaluation, and the behaviour
of critical eigenfunctions. In particular, a harmonic identity for
a candidate critical first-promotion kernel can require ruling out
escape without a future core promotion. That critical-kernel issue
is not an assumption in the computability proof above: the upper
certificates use submultiplicative finite-time means at discounts
strictly above $\mu$.

\paragraph{Verification.}
The analytical results are mathematical proofs supported by exact
finite certificates. The infection reproduction bound is checked by
\texttt{scripts/verify\_infection\_kernel.py}; the beta burst and
generation calculations are checked by
\texttt{scripts/core\_clock\_kernel.py} and
\texttt{scripts/inert\_growth\_kernel.py}. The numerical refinement
has the separate full check described in Section~\ref{numerical:section}.
These computations verify finite arithmetic claims and identities;
they do not replace the infinite-state coupling and termination
arguments. This research draft is not a proof-assistant formalization
or a claim of external peer review.

\paragraph{Acknowledgments and provenance.}
The work builds on Krapivsky and Redner's model and Uzay's exponent
existence and computer-assisted bounds. The motif tables and their
original drift decomposition are attributed to the pinned upstream
source. AI systems assisted with derivations, proof audits, and
verification code during preparation; responsibility for the
mathematical statements rests with the author. The repository records
the imported inputs and reproducible finite checks.


\begin{thebibliography}{9}

\bibitem{KR}
P. L. Krapivsky and S. Redner,
\emph{Emergent Network Modularity},
Journal of Statistical Mechanics: Theory and Experiment (2017), 073405.
\href{https://doi.org/10.1088/1742-5468/aa7a3f}{doi:10.1088/1742-5468/aa7a3f}.
\href{https://arxiv.org/abs/1706.01514}{arXiv:1706.01514}.

\bibitem{AB}
L. Addario-Berry, S. Briend, L. Devroye, S. Donderwinkel,
C. Kerriou, and G. Lugosi,
\emph{Random Friend Trees}, Random Structures \& Algorithms
\textbf{69} (2026), e70082.
\href{https://doi.org/10.1002/rsa.70082}{doi:10.1002/rsa.70082}.
\href{https://arxiv.org/abs/2403.20185}{arXiv:2403.20185}.

\bibitem{Uzay}
Can Uzay,
\emph{The random friend tree exponent is not 17/30},
proofs and exact-arithmetic verification repository, 2026.
Existence and planted-branch input: \texttt{proofs/existence.md},
especially Lemma 4.1(c), Corollary 5.3, and Lemma 6.1;
almost-sure identification: \texttt{proofs/almost-sure-exponent.md}.
Revision
\href{https://github.com/uzay-eth/friend-tree-exponent/tree/88f2f9c8c3ebc1b4e0c980db6afd97d29f3df176}{\texttt{88f2f9c8c3ebc1b4e0c980db6afd97d29f3df176}}.

\bibitem{UzayTables}
Can Uzay, upper motif certificate in
\emph{The random friend tree exponent is not 17/30}, 2026,
\texttt{verify/motif-bracket/}.
Numerical inputs pinned at revision
\href{https://github.com/uzay-eth/friend-tree-exponent/tree/ba2ec2e1935e0858e6bae2b40a9ee07c9b5665aa}{\texttt{ba2ec2e1935e0858e6bae2b40a9ee07c9b5665aa}}.
The exact source and binary hashes are recorded in the companion
repository's upstream manifest.

\bibitem{Local}
C. D. Long,
\emph{Growth and irrationality in local redirection trees},
companion research draft, November 23, 2026.
Source and PDF: \texttt{local\_redirection\_irrationality.tex} and
\texttt{local\_redirection\_irrationality.pdf} in the same repository.

\end{thebibliography}
\end{document}
